% Auto-generated by scripts/extract_catalog.py — do not edit.

\section*{\texttt{theories/foundations/prelude.v}}
\begin{description}\setlength\itemsep{0pt}\small
\item[\texttt{prelude\_classical}] (lemma) \href{https://github.com/LLM4Rocq/digraph-theory/blob/main/theories/foundations/prelude.v\#L29}{\texttt{prelude.v:29}} --- Classical logic is in scope (derived from \texttt{boolp}'s axioms).
\item[\texttt{prelude\_set}] (lemma) \href{https://github.com/LLM4Rocq/digraph-theory/blob/main/theories/foundations/prelude.v\#L33}{\texttt{prelude.v:33}} --- The classical set layer is in scope.
\item[\texttt{prelude\_ssr}] (lemma) \href{https://github.com/LLM4Rocq/digraph-theory/blob/main/theories/foundations/prelude.v\#L37}{\texttt{prelude.v:37}} --- MathComp's small-scale reflection is in scope.
\item[\texttt{card\_classes}] (lemma) \href{https://github.com/LLM4Rocq/digraph-theory/blob/main/theories/foundations/prelude.v\#L47}{\texttt{prelude.v:47}} --- ** General counting Partition counting by a bounded nat-valued class function: any finite set splits as the sum of its class sizes.
\item[\texttt{card\_classes\_inj}] (lemma) \href{https://github.com/LLM4Rocq/digraph-theory/blob/main/theories/foundations/prelude.v\#L68}{\texttt{prelude.v:68}} --- When the class function is injective on \texttt{K} (one element per class at most --- e.g.
\end{description}

\section*{\texttt{theories/foundations/interop\_graph\_theory.v}}
\begin{description}\setlength\itemsep{0pt}\small
\item[\texttt{clique\_set1}] (lemma) \href{https://github.com/LLM4Rocq/digraph-theory/blob/main/theories/foundations/interop_graph_theory.v\#L30}{\texttt{interop\_graph\_theory.v:30}}
\item[\texttt{set1\_cliques}] (lemma) \href{https://github.com/LLM4Rocq/digraph-theory/blob/main/theories/foundations/interop_graph_theory.v\#L33}{\texttt{interop\_graph\_theory.v:33}}
\item[\texttt{omega\_ge1}] (lemma) \href{https://github.com/LLM4Rocq/digraph-theory/blob/main/theories/foundations/interop_graph_theory.v\#L36}{\texttt{interop\_graph\_theory.v:36}}
\item[\texttt{omega\_le\_card}] (lemma) \href{https://github.com/LLM4Rocq/digraph-theory/blob/main/theories/foundations/interop_graph_theory.v\#L41}{\texttt{interop\_graph\_theory.v:41}}
\item[\texttt{omega\_le1}] (lemma) \href{https://github.com/LLM4Rocq/digraph-theory/blob/main/theories/foundations/interop_graph_theory.v\#L47}{\texttt{interop\_graph\_theory.v:47}} --- An edgeless subset has clique number at most \texttt{1}.
\item[\texttt{omega\_ge2}] (lemma) \href{https://github.com/LLM4Rocq/digraph-theory/blob/main/theories/foundations/interop_graph_theory.v\#L56}{\texttt{interop\_graph\_theory.v:56}} --- Conversely, a single edge pushes $\omega$ to at least 2.
\item[\texttt{omega\_hom}] (lemma) \href{https://github.com/LLM4Rocq/digraph-theory/blob/main/theories/foundations/interop_graph_theory.v\#L69}{\texttt{interop\_graph\_theory.v:69}} --- $\omega$ is monotone along injective edge-preserving maps (used to compare backedge graphs of a sub-tournament and its host).
\item[\texttt{K2\_rel}] (definition) \href{https://github.com/LLM4Rocq/digraph-theory/blob/main/theories/foundations/interop_graph_theory.v\#L88}{\texttt{interop\_graph\_theory.v:88}} --- ** Smoke check --- $\omega$ of a tiny hand-built \texttt{sgraph} The original M0 exit criterion: graph-theory's clique number is usable on a concrete graph.
\item[\texttt{K2\_sym}] (fact) \href{https://github.com/LLM4Rocq/digraph-theory/blob/main/theories/foundations/interop_graph_theory.v\#L89}{\texttt{interop\_graph\_theory.v:89}}
\item[\texttt{K2\_irrefl}] (fact) \href{https://github.com/LLM4Rocq/digraph-theory/blob/main/theories/foundations/interop_graph_theory.v\#L91}{\texttt{interop\_graph\_theory.v:91}}
\item[\texttt{K2}] (definition) \href{https://github.com/LLM4Rocq/digraph-theory/blob/main/theories/foundations/interop_graph_theory.v\#L93}{\texttt{interop\_graph\_theory.v:93}}
\item[\texttt{omega\_K2}] (lemma) \href{https://github.com/LLM4Rocq/digraph-theory/blob/main/theories/foundations/interop_graph_theory.v\#L95}{\texttt{interop\_graph\_theory.v:95}}
\end{description}

\section*{\texttt{theories/core/digraph.v}}
\begin{description}\setlength\itemsep{0pt}\small
\item[\texttt{to\_GT}] (definition) \href{https://github.com/LLM4Rocq/digraph-theory/blob/main/theories/core/digraph.v\#L56}{\texttt{digraph.v:56}} --- ** Projection into graph-theory's record world (D7 glue) \texttt{diGraph} is graph-theory's notation for its \texttt{relType} record (re-exported by the interop file); \texttt{RelType} is its constructor --- its \texttt{DiGraph} alias is shadowed by our structure module above.
\item[\texttt{converse}] (definition) \href{https://github.com/LLM4Rocq/digraph-theory/blob/main/theories/core/digraph.v\#L60}{\texttt{digraph.v:60}} --- ** Converse digraph (type alias, mathcomp \texttt{T\textasciicircum{}d}-style)
\item[\texttt{converse\_arcE}] (lemma) \href{https://github.com/LLM4Rocq/digraph-theory/blob/main/theories/core/digraph.v\#L66}{\texttt{digraph.v:66}}
\item[\texttt{sub\_arcE}] (lemma) \href{https://github.com/LLM4Rocq/digraph-theory/blob/main/theories/core/digraph.v\#L78}{\texttt{digraph.v:78}}
\item[\texttt{induced\_digraph}] (definition) \href{https://github.com/LLM4Rocq/digraph-theory/blob/main/theories/core/digraph.v\#L86}{\texttt{digraph.v:86}} --- Object-level versions, convenient in statements (the \texttt{: diGraphType} cast goes through Rocq's reverse-coercion via the canonical instances).
\item[\texttt{del\_vertex}] (definition) \href{https://github.com/LLM4Rocq/digraph-theory/blob/main/theories/core/digraph.v\#L89}{\texttt{digraph.v:89}}
\item[\texttt{dgiso}] (definition) \href{https://github.com/LLM4Rocq/digraph-theory/blob/main/theories/core/digraph.v\#L97}{\texttt{digraph.v:97}} --- ** Digraph isomorphism (Named \texttt{dgiso}: graph-theory's \texttt{diso} --- re-exported through the interop --- is the *simple-graph* iso
\item[\texttt{dgiso\_refl}] (lemma) \href{https://github.com/LLM4Rocq/digraph-theory/blob/main/theories/core/digraph.v\#L100}{\texttt{digraph.v:100}}
\item[\texttt{dgiso\_sym}] (lemma) \href{https://github.com/LLM4Rocq/digraph-theory/blob/main/theories/core/digraph.v\#L103}{\texttt{digraph.v:103}}
\item[\texttt{dgiso\_trans}] (lemma) \href{https://github.com/LLM4Rocq/digraph-theory/blob/main/theories/core/digraph.v\#L109}{\texttt{digraph.v:109}}
\item[\texttt{dgiso\_card}] (lemma) \href{https://github.com/LLM4Rocq/digraph-theory/blob/main/theories/core/digraph.v\#L117}{\texttt{digraph.v:117}}
\end{description}

\section*{\texttt{theories/core/oriented.v}}
\begin{description}\setlength\itemsep{0pt}\small
\item[\texttt{outdeg}] (definition) \href{https://github.com/LLM4Rocq/digraph-theory/blob/main/theories/core/oriented.v\#L47}{\texttt{oriented.v:47}}
\item[\texttt{outdeg\_in}] (definition) \href{https://github.com/LLM4Rocq/digraph-theory/blob/main/theories/core/oriented.v\#L51}{\texttt{oriented.v:51}} --- Out-degree into a fixed vertex subset --- the induced out-degree when \texttt{v \textbackslash{}in A}.
\item[\texttt{outdeg\_inT}] (lemma) \href{https://github.com/LLM4Rocq/digraph-theory/blob/main/theories/core/oriented.v\#L53}{\texttt{oriented.v:53}}
\item[\texttt{outdeg\_in\_le}] (lemma) \href{https://github.com/LLM4Rocq/digraph-theory/blob/main/theories/core/oriented.v\#L56}{\texttt{oriented.v:56}}
\item[\texttt{outdeg\_in\_mono}] (lemma) \href{https://github.com/LLM4Rocq/digraph-theory/blob/main/theories/core/oriented.v\#L61}{\texttt{oriented.v:61}}
\item[\texttt{outdeg\_in\_sumE}] (lemma) \href{https://github.com/LLM4Rocq/digraph-theory/blob/main/theories/core/oriented.v\#L67}{\texttt{oriented.v:67}}
\item[\texttt{outdeg\_induced}] (lemma) \href{https://github.com/LLM4Rocq/digraph-theory/blob/main/theories/core/oriented.v\#L76}{\texttt{oriented.v:76}} --- The induced subgraph on \texttt{A} has the inner out-degrees.
\item[\texttt{arc\_pair\_bound}] (lemma) \href{https://github.com/LLM4Rocq/digraph-theory/blob/main/theories/core/oriented.v\#L92}{\texttt{oriented.v:92}}
\item[\texttt{oriented\_arcs\_bound}] (lemma) \href{https://github.com/LLM4Rocq/digraph-theory/blob/main/theories/core/oriented.v\#L100}{\texttt{oriented.v:100}} --- O1, counting form: twice the arc count of D\texttt{A} is $\le$ |A|(|A|--1).
\item[\texttt{oriented\_avg\_bound}] (lemma) \href{https://github.com/LLM4Rocq/digraph-theory/blob/main/theories/core/oriented.v\#L119}{\texttt{oriented.v:119}} --- O1(a): some vertex of a nonempty A has inner out-degree $\le$ (|A|--1)./2.
\item[\texttt{oriented\_mindeg\_card}] (lemma) \href{https://github.com/LLM4Rocq/digraph-theory/blob/main/theories/core/oriented.v\#L144}{\texttt{oriented.v:144}} --- O1(b): if all inner out-degrees are at least k, then |A| $\ge$ 2k+1.
\item[\texttt{oriented\_card}] (lemma) \href{https://github.com/LLM4Rocq/digraph-theory/blob/main/theories/core/oriented.v\#L159}{\texttt{oriented.v:159}} --- O1(c): a nonempty oriented digraph with min out-degree $\ge$ k has $\ge$ 2k+1 vertices.
\item[\texttt{sub\_oarc\_irrefl}] (fact) \href{https://github.com/LLM4Rocq/digraph-theory/blob/main/theories/core/oriented.v\#L175}{\texttt{oriented.v:175}}
\item[\texttt{sub\_oarc\_asymm}] (fact) \href{https://github.com/LLM4Rocq/digraph-theory/blob/main/theories/core/oriented.v\#L178}{\texttt{oriented.v:178}}
\item[\texttt{induced\_oriented}] (definition) \href{https://github.com/LLM4Rocq/digraph-theory/blob/main/theories/core/oriented.v\#L187}{\texttt{oriented.v:187}} --- Object-level version for statements.
\item[\texttt{outsel}] (definition) \href{https://github.com/LLM4Rocq/digraph-theory/blob/main/theories/core/oriented.v\#L195}{\texttt{oriented.v:195}} --- The alias takes \texttt{f} as a REAL parameter (M2 lesson: a type alias must use its parameters, or section discharge drops them and all selections collapse onto one type).
\item[\texttt{outsel\_arcE}] (lemma) \href{https://github.com/LLM4Rocq/digraph-theory/blob/main/theories/core/oriented.v\#L204}{\texttt{oriented.v:204}}
\item[\texttt{outsel\_arc\_sub}] (lemma) \href{https://github.com/LLM4Rocq/digraph-theory/blob/main/theories/core/oriented.v\#L208}{\texttt{oriented.v:208}}
\item[\texttt{outsel\_irrefl}] (fact) \href{https://github.com/LLM4Rocq/digraph-theory/blob/main/theories/core/oriented.v\#L219}{\texttt{oriented.v:219}}
\item[\texttt{outsel\_asymm}] (fact) \href{https://github.com/LLM4Rocq/digraph-theory/blob/main/theories/core/oriented.v\#L222}{\texttt{oriented.v:222}}
\item[\texttt{ksel}] (definition) \href{https://github.com/LLM4Rocq/digraph-theory/blob/main/theories/core/oriented.v\#L237}{\texttt{oriented.v:237}}
\item[\texttt{ksel\_sub}] (lemma) \href{https://github.com/LLM4Rocq/digraph-theory/blob/main/theories/core/oriented.v\#L239}{\texttt{oriented.v:239}}
\item[\texttt{ksel\_card}] (lemma) \href{https://github.com/LLM4Rocq/digraph-theory/blob/main/theories/core/oriented.v\#L245}{\texttt{oriented.v:245}}
\item[\texttt{outsel\_ksel\_outdeg}] (lemma) \href{https://github.com/LLM4Rocq/digraph-theory/blob/main/theories/core/oriented.v\#L254}{\texttt{oriented.v:254}} --- O2's exit: in \texttt{outsel ksel}, every out-degree is exactly k.
\end{description}

\section*{\texttt{theories/core/dipath.v}}
\begin{description}\setlength\itemsep{0pt}\small
\item[\texttt{dipath}] (definition) \href{https://github.com/LLM4Rocq/digraph-theory/blob/main/theories/core/dipath.v\#L36}{\texttt{dipath.v:36}} --- ** Directed simple paths
\item[\texttt{dipath\_path}] (lemma) \href{https://github.com/LLM4Rocq/digraph-theory/blob/main/theories/core/dipath.v\#L38}{\texttt{dipath.v:38}}
\item[\texttt{dipath\_uniq}] (lemma) \href{https://github.com/LLM4Rocq/digraph-theory/blob/main/theories/core/dipath.v\#L41}{\texttt{dipath.v:41}}
\item[\texttt{dipath\_nil}] (lemma) \href{https://github.com/LLM4Rocq/digraph-theory/blob/main/theories/core/dipath.v\#L44}{\texttt{dipath.v:44}}
\item[\texttt{dipath\_rcons}] (lemma) \href{https://github.com/LLM4Rocq/digraph-theory/blob/main/theories/core/dipath.v\#L47}{\texttt{dipath.v:47}}
\item[\texttt{cat\_dipath}] (lemma) \href{https://github.com/LLM4Rocq/digraph-theory/blob/main/theories/core/dipath.v\#L56}{\texttt{dipath.v:56}}
\item[\texttt{cat\_dipath\_cont}] (lemma) \href{https://github.com/LLM4Rocq/digraph-theory/blob/main/theories/core/dipath.v\#L67}{\texttt{dipath.v:67}} --- Continuation form: append a path that starts AT the endpoint.
\item[\texttt{dipath\_cons}] (lemma) \href{https://github.com/LLM4Rocq/digraph-theory/blob/main/theories/core/dipath.v\#L83}{\texttt{dipath.v:83}}
\item[\texttt{dipath\_drop}] (lemma) \href{https://github.com/LLM4Rocq/digraph-theory/blob/main/theories/core/dipath.v\#L89}{\texttt{dipath.v:89}}
\item[\texttt{last\_take}] (lemma) \href{https://github.com/LLM4Rocq/digraph-theory/blob/main/theories/core/dipath.v\#L105}{\texttt{dipath.v:105}}
\item[\texttt{last\_drop}] (lemma) \href{https://github.com/LLM4Rocq/digraph-theory/blob/main/theories/core/dipath.v\#L108}{\texttt{dipath.v:108}}
\item[\texttt{take\_path\_last}] (lemma) \href{https://github.com/LLM4Rocq/digraph-theory/blob/main/theories/core/dipath.v\#L117}{\texttt{dipath.v:117}} --- Endpoint of a prefix, on the (x :: s) indexing.
\item[\texttt{dipath\_take}] (lemma) \href{https://github.com/LLM4Rocq/digraph-theory/blob/main/theories/core/dipath.v\#L124}{\texttt{dipath.v:124}}
\item[\texttt{dipath\_arc\_nth}] (lemma) \href{https://github.com/LLM4Rocq/digraph-theory/blob/main/theories/core/dipath.v\#L132}{\texttt{dipath.v:132}} --- The i-th vertex of the path (x, s), i ranging over 0..size s.
\item[\texttt{dipath\_size}] (lemma) \href{https://github.com/LLM4Rocq/digraph-theory/blob/main/theories/core/dipath.v\#L136}{\texttt{dipath.v:136}}
\item[\texttt{has\_dipath}] (definition) \href{https://github.com/LLM4Rocq/digraph-theory/blob/main/theories/core/dipath.v\#L144}{\texttt{dipath.v:144}} --- ** The longest-path length $\ell$(D)
\item[\texttt{has\_dipathP}] (lemma) \href{https://github.com/LLM4Rocq/digraph-theory/blob/main/theories/core/dipath.v\#L146}{\texttt{dipath.v:146}}
\item[\texttt{ell}] (definition) \href{https://github.com/LLM4Rocq/digraph-theory/blob/main/theories/core/dipath.v\#L155}{\texttt{dipath.v:155}}
\item[\texttt{ell\_max}] (lemma) \href{https://github.com/LLM4Rocq/digraph-theory/blob/main/theories/core/dipath.v\#L157}{\texttt{dipath.v:157}}
\item[\texttt{ellP}] (lemma) \href{https://github.com/LLM4Rocq/digraph-theory/blob/main/theories/core/dipath.v\#L166}{\texttt{dipath.v:166}}
\item[\texttt{maxpath\_endclosure}] (lemma) \href{https://github.com/LLM4Rocq/digraph-theory/blob/main/theories/core/dipath.v\#L178}{\texttt{dipath.v:178}} --- ** K-A: out-neighbours of the endpoint of a maximum path lie on it
\item[\texttt{endmax}] (definition) \href{https://github.com/LLM4Rocq/digraph-theory/blob/main/theories/core/dipath.v\#L188}{\texttt{dipath.v:188}} --- ** End-maximal paths and K-Ext
\item[\texttt{endmax\_closure}] (lemma) \href{https://github.com/LLM4Rocq/digraph-theory/blob/main/theories/core/dipath.v\#L191}{\texttt{dipath.v:191}}
\item[\texttt{endmax\_dipath}] (lemma) \href{https://github.com/LLM4Rocq/digraph-theory/blob/main/theories/core/dipath.v\#L195}{\texttt{dipath.v:195}}
\item[\texttt{endmax\_ex}] (lemma) \href{https://github.com/LLM4Rocq/digraph-theory/blob/main/theories/core/dipath.v\#L199}{\texttt{dipath.v:199}} --- K-Ext: every vertex starts some end-maximal path.
\item[\texttt{dicycle}] (definition) \href{https://github.com/LLM4Rocq/digraph-theory/blob/main/theories/core/dipath.v\#L223}{\texttt{dipath.v:223}} --- ** Directed cycles as seqs
\item[\texttt{dicycle\_rot}] (lemma) \href{https://github.com/LLM4Rocq/digraph-theory/blob/main/theories/core/dipath.v\#L225}{\texttt{dipath.v:225}}
\item[\texttt{dicycle\_next}] (lemma) \href{https://github.com/LLM4Rocq/digraph-theory/blob/main/theories/core/dipath.v\#L229}{\texttt{dipath.v:229}} --- The closing/cyclic arcs: every cycle vertex beats its \texttt{next}.
\item[\texttt{dicycle\_prev}] (lemma) \href{https://github.com/LLM4Rocq/digraph-theory/blob/main/theories/core/dipath.v\#L232}{\texttt{dipath.v:232}}
\item[\texttt{dicycle\_unroll}] (lemma) \href{https://github.com/LLM4Rocq/digraph-theory/blob/main/theories/core/dipath.v\#L237}{\texttt{dipath.v:237}} --- Unrolling: a cycle yields a path from any of its vertices through all of it, ending at the predecessor (with the closing arc back).
\item[\texttt{dicycle\_suffix}] (lemma) \href{https://github.com/LLM4Rocq/digraph-theory/blob/main/theories/core/dipath.v\#L262}{\texttt{dipath.v:262}} --- Cycle from a path suffix plus a back-arc from the endpoint.
\item[\texttt{prev\_head}] (lemma) \href{https://github.com/LLM4Rocq/digraph-theory/blob/main/theories/core/dipath.v\#L290}{\texttt{dipath.v:290}} --- \texttt{prev} of the head of a uniq seq is its last element (eqType-level helper for cycle predecessor reasoning).
\item[\texttt{dipath\_map}] (lemma) \href{https://github.com/LLM4Rocq/digraph-theory/blob/main/theories/core/dipath.v\#L304}{\texttt{dipath.v:304}} --- ** M-$\ell$: $\ell$ is monotone under injective arc-preserving embeddings
\item[\texttt{ell\_embed}] (lemma) \href{https://github.com/LLM4Rocq/digraph-theory/blob/main/theories/core/dipath.v\#L314}{\texttt{dipath.v:314}}
\item[\texttt{ell\_outsel}] (lemma) \href{https://github.com/LLM4Rocq/digraph-theory/blob/main/theories/core/dipath.v\#L330}{\texttt{dipath.v:330}} --- The two instances used by the reduction R: arc-sub-relations and induced subgraphs.
\item[\texttt{ell\_induced}] (lemma) \href{https://github.com/LLM4Rocq/digraph-theory/blob/main/theories/core/dipath.v\#L337}{\texttt{dipath.v:337}}
\end{description}

\section*{\texttt{theories/core/tournament.v}}
\begin{description}\setlength\itemsep{0pt}\small
\item[\texttt{asymm}] (fact) \href{https://github.com/LLM4Rocq/digraph-theory/blob/main/theories/core/tournament.v\#L47}{\texttt{tournament.v:47}}
\item[\texttt{arcxx}] (lemma) \href{https://github.com/LLM4Rocq/digraph-theory/blob/main/theories/core/tournament.v\#L65}{\texttt{tournament.v:65}}
\item[\texttt{arc\_neq}] (lemma) \href{https://github.com/LLM4Rocq/digraph-theory/blob/main/theories/core/tournament.v\#L68}{\texttt{tournament.v:68}}
\item[\texttt{arc\_asym}] (lemma) \href{https://github.com/LLM4Rocq/digraph-theory/blob/main/theories/core/tournament.v\#L71}{\texttt{tournament.v:71}}
\item[\texttt{arcNarc}] (lemma) \href{https://github.com/LLM4Rocq/digraph-theory/blob/main/theories/core/tournament.v\#L78}{\texttt{tournament.v:78}} --- On distinct vertices, the two arc directions are complementary.
\item[\texttt{arc\_or}] (lemma) \href{https://github.com/LLM4Rocq/digraph-theory/blob/main/theories/core/tournament.v\#L83}{\texttt{tournament.v:83}}
\item[\texttt{arc\_xorE}] (lemma) \href{https://github.com/LLM4Rocq/digraph-theory/blob/main/theories/core/tournament.v\#L87}{\texttt{tournament.v:87}} --- Since both directions never hold together, xor and or coincide.
\item[\texttt{N\_out}] (definition) \href{https://github.com/LLM4Rocq/digraph-theory/blob/main/theories/core/tournament.v\#L95}{\texttt{tournament.v:95}} --- ** Neighbourhoods
\item[\texttt{N\_in}] (definition) \href{https://github.com/LLM4Rocq/digraph-theory/blob/main/theories/core/tournament.v\#L96}{\texttt{tournament.v:96}}
\item[\texttt{N\_outE}] (lemma) \href{https://github.com/LLM4Rocq/digraph-theory/blob/main/theories/core/tournament.v\#L98}{\texttt{tournament.v:98}}
\item[\texttt{N\_inE}] (lemma) \href{https://github.com/LLM4Rocq/digraph-theory/blob/main/theories/core/tournament.v\#L99}{\texttt{tournament.v:99}}
\item[\texttt{N\_out\_in\_partition}] (lemma) \href{https://github.com/LLM4Rocq/digraph-theory/blob/main/theories/core/tournament.v\#L101}{\texttt{tournament.v:101}}
\item[\texttt{transb}] (definition) \href{https://github.com/LLM4Rocq/digraph-theory/blob/main/theories/core/tournament.v\#L112}{\texttt{tournament.v:112}} --- ** Decidable transitivity and the 3-cycle dichotomy
\item[\texttt{transbP}] (lemma) \href{https://github.com/LLM4Rocq/digraph-theory/blob/main/theories/core/tournament.v\#L115}{\texttt{tournament.v:115}}
\item[\texttt{ntransbP}] (lemma) \href{https://github.com/LLM4Rocq/digraph-theory/blob/main/theories/core/tournament.v\#L124}{\texttt{tournament.v:124}} --- A tournament fails transitivity exactly when it has a directed triangle.
\item[\texttt{ntransb\_card}] (lemma) \href{https://github.com/LLM4Rocq/digraph-theory/blob/main/theories/core/tournament.v\#L139}{\texttt{tournament.v:139}} --- An intransitive tournament has at least 3 vertices.
\item[\texttt{card\_le2\_transb}] (lemma) \href{https://github.com/LLM4Rocq/digraph-theory/blob/main/theories/core/tournament.v\#L148}{\texttt{tournament.v:148}}
\item[\texttt{sub\_arc\_total}] (fact) \href{https://github.com/LLM4Rocq/digraph-theory/blob/main/theories/core/tournament.v\#L160}{\texttt{tournament.v:160}} --- The subtype already carries the \texttt{Oriented} structure (core/oriented.v); only totality is added here.
\item[\texttt{sub\_tournament}] (definition) \href{https://github.com/LLM4Rocq/digraph-theory/blob/main/theories/core/tournament.v\#L171}{\texttt{tournament.v:171}} --- Object-level versions for use in statements.
\item[\texttt{del\_tournament}] (definition) \href{https://github.com/LLM4Rocq/digraph-theory/blob/main/theories/core/tournament.v\#L174}{\texttt{tournament.v:174}}
\item[\texttt{C3}] (definition) \href{https://github.com/LLM4Rocq/digraph-theory/blob/main/theories/core/tournament.v\#L183}{\texttt{tournament.v:183}}
\item[\texttt{C3\_irrefl}] (fact) \href{https://github.com/LLM4Rocq/digraph-theory/blob/main/theories/core/tournament.v\#L188}{\texttt{tournament.v:188}}
\item[\texttt{C3\_total}] (fact) \href{https://github.com/LLM4Rocq/digraph-theory/blob/main/theories/core/tournament.v\#L191}{\texttt{tournament.v:191}}
\item[\texttt{arcC3E}] (lemma) \href{https://github.com/LLM4Rocq/digraph-theory/blob/main/theories/core/tournament.v\#L196}{\texttt{tournament.v:196}}
\item[\texttt{card\_C3}] (lemma) \href{https://github.com/LLM4Rocq/digraph-theory/blob/main/theories/core/tournament.v\#L199}{\texttt{tournament.v:199}}
\item[\texttt{C3\_Ntransb}] (lemma) \href{https://github.com/LLM4Rocq/digraph-theory/blob/main/theories/core/tournament.v\#L203}{\texttt{tournament.v:203}} --- C3 is the directed triangle: it is *not* transitive.
\item[\texttt{TT}] (definition) \href{https://github.com/LLM4Rocq/digraph-theory/blob/main/theories/core/tournament.v\#L215}{\texttt{tournament.v:215}}
\item[\texttt{TT\_irrefl}] (fact) \href{https://github.com/LLM4Rocq/digraph-theory/blob/main/theories/core/tournament.v\#L220}{\texttt{tournament.v:220}}
\item[\texttt{TT\_total}] (fact) \href{https://github.com/LLM4Rocq/digraph-theory/blob/main/theories/core/tournament.v\#L223}{\texttt{tournament.v:223}}
\item[\texttt{arcTTE}] (lemma) \href{https://github.com/LLM4Rocq/digraph-theory/blob/main/theories/core/tournament.v\#L232}{\texttt{tournament.v:232}}
\item[\texttt{card\_TT}] (lemma) \href{https://github.com/LLM4Rocq/digraph-theory/blob/main/theories/core/tournament.v\#L235}{\texttt{tournament.v:235}}
\item[\texttt{TT\_transb}] (lemma) \href{https://github.com/LLM4Rocq/digraph-theory/blob/main/theories/core/tournament.v\#L238}{\texttt{tournament.v:238}}
\end{description}

\section*{\texttt{theories/core/order.v}}
\begin{description}\setlength\itemsep{0pt}\small
\item[\texttt{ltp}] (definition) \href{https://github.com/LLM4Rocq/digraph-theory/blob/main/theories/core/order.v\#L39}{\texttt{order.v:39}}
\item[\texttt{ltp\_irrefl}] (lemma) \href{https://github.com/LLM4Rocq/digraph-theory/blob/main/theories/core/order.v\#L41}{\texttt{order.v:41}}
\item[\texttt{ltp\_trans}] (lemma) \href{https://github.com/LLM4Rocq/digraph-theory/blob/main/theories/core/order.v\#L44}{\texttt{order.v:44}}
\item[\texttt{ltp\_total}] (lemma) \href{https://github.com/LLM4Rocq/digraph-theory/blob/main/theories/core/order.v\#L47}{\texttt{order.v:47}}
\item[\texttt{ltp\_asym}] (lemma) \href{https://github.com/LLM4Rocq/digraph-theory/blob/main/theories/core/order.v\#L53}{\texttt{order.v:53}}
\item[\texttt{le\_r\_total}] (fact) \href{https://github.com/LLM4Rocq/digraph-theory/blob/main/theories/core/order.v\#L66}{\texttt{order.v:66}}
\item[\texttt{le\_r\_trans}] (fact) \href{https://github.com/LLM4Rocq/digraph-theory/blob/main/theories/core/order.v\#L73}{\texttt{order.v:73}}
\item[\texttt{mem\_s}] (fact) \href{https://github.com/LLM4Rocq/digraph-theory/blob/main/theories/core/order.v\#L82}{\texttt{order.v:82}}
\item[\texttt{size\_s}] (fact) \href{https://github.com/LLM4Rocq/digraph-theory/blob/main/theories/core/order.v\#L83}{\texttt{order.v:83}}
\item[\texttt{sorted\_s}] (fact) \href{https://github.com/LLM4Rocq/digraph-theory/blob/main/theories/core/order.v\#L84}{\texttt{order.v:84}}
\item[\texttt{rank\_lt}] (fact) \href{https://github.com/LLM4Rocq/digraph-theory/blob/main/theories/core/order.v\#L85}{\texttt{order.v:85}}
\item[\texttt{rfun}] (definition) \href{https://github.com/LLM4Rocq/digraph-theory/blob/main/theories/core/order.v\#L87}{\texttt{order.v:87}}
\item[\texttt{rfun\_inj}] (fact) \href{https://github.com/LLM4Rocq/digraph-theory/blob/main/theories/core/order.v\#L89}{\texttt{order.v:89}}
\item[\texttt{realize}] (definition) \href{https://github.com/LLM4Rocq/digraph-theory/blob/main/theories/core/order.v\#L96}{\texttt{order.v:96}}
\item[\texttt{ltp\_realizeE}] (lemma) \href{https://github.com/LLM4Rocq/digraph-theory/blob/main/theories/core/order.v\#L98}{\texttt{order.v:98}}
\item[\texttt{ltp\_pullback}] (lemma) \href{https://github.com/LLM4Rocq/digraph-theory/blob/main/theories/core/order.v\#L119}{\texttt{order.v:119}} --- Transport an order along an injective map: the pulled-back order on \texttt{U} is again realized by a permutation.
\item[\texttt{backedge\_rel}] (definition) \href{https://github.com/LLM4Rocq/digraph-theory/blob/main/theories/core/order.v\#L136}{\texttt{order.v:136}}
\item[\texttt{backedge\_sym}] (fact) \href{https://github.com/LLM4Rocq/digraph-theory/blob/main/theories/core/order.v\#L139}{\texttt{order.v:139}}
\item[\texttt{backedge\_irrefl}] (fact) \href{https://github.com/LLM4Rocq/digraph-theory/blob/main/theories/core/order.v\#L142}{\texttt{order.v:142}}
\item[\texttt{backedge}] (definition) \href{https://github.com/LLM4Rocq/digraph-theory/blob/main/theories/core/order.v\#L145}{\texttt{order.v:145}}
\item[\texttt{backedgeE}] (lemma) \href{https://github.com/LLM4Rocq/digraph-theory/blob/main/theories/core/order.v\#L147}{\texttt{order.v:147}}
\item[\texttt{backedge\_neq}] (lemma) \href{https://github.com/LLM4Rocq/digraph-theory/blob/main/theories/core/order.v\#L153}{\texttt{order.v:153}} --- An edge of the backedge graph is in particular an arc-disagreement, so its endpoints are distinct.
\item[\texttt{backedge\_arc\_forward}] (lemma) \href{https://github.com/LLM4Rocq/digraph-theory/blob/main/theories/core/order.v\#L163}{\texttt{order.v:163}} --- If the order linearizes the arcs (no backward arc at all), the backedge graph is edgeless.
\end{description}

\section*{\texttt{theories/invariants/omegabar.v}}
\begin{description}\setlength\itemsep{0pt}\small
\item[\texttt{omegab\_at}] (definition) \href{https://github.com/LLM4Rocq/digraph-theory/blob/main/theories/invariants/omegabar.v\#L29}{\texttt{omegabar.v:29}} --- ** Definition
\item[\texttt{omegabar}] (definition) \href{https://github.com/LLM4Rocq/digraph-theory/blob/main/theories/invariants/omegabar.v\#L32}{\texttt{omegabar.v:32}}
\item[\texttt{omegabar\_min}] (lemma) \href{https://github.com/LLM4Rocq/digraph-theory/blob/main/theories/invariants/omegabar.v\#L42}{\texttt{omegabar.v:42}}
\item[\texttt{omegabar\_witness}] (lemma) \href{https://github.com/LLM4Rocq/digraph-theory/blob/main/theories/invariants/omegabar.v\#L45}{\texttt{omegabar.v:45}}
\item[\texttt{omegabar\_gt0}] (lemma) \href{https://github.com/LLM4Rocq/digraph-theory/blob/main/theories/invariants/omegabar.v\#L48}{\texttt{omegabar.v:48}}
\item[\texttt{omegab\_at\_le1}] (lemma) \href{https://github.com/LLM4Rocq/digraph-theory/blob/main/theories/invariants/omegabar.v\#L56}{\texttt{omegabar.v:56}} --- ** $\bar\omega$ = 1 characterizes transitivity
\item[\texttt{omegabar\_le1}] (lemma) \href{https://github.com/LLM4Rocq/digraph-theory/blob/main/theories/invariants/omegabar.v\#L63}{\texttt{omegabar.v:63}}
\item[$\bigstar$ \texttt{omegabar\_transb}] (lemma) \href{https://github.com/LLM4Rocq/digraph-theory/blob/main/theories/invariants/omegabar.v\#L74}{\texttt{omegabar.v:74}}
\item[\texttt{omegabar\_embed}] (lemma) \href{https://github.com/LLM4Rocq/digraph-theory/blob/main/theories/invariants/omegabar.v\#L108}{\texttt{omegabar.v:108}} --- ** Monotonicity under arc-preserving embeddings Any injective map that reflects arcs both ways pushes $\bar\omega$ up: the witness order on the big tournament pulls back along the embedding (\texttt{ltp\_pullback}) and the backedge graphs correspond edge-for-edge.
\item[\texttt{omegabar\_dgiso}] (lemma) \href{https://github.com/LLM4Rocq/digraph-theory/blob/main/theories/invariants/omegabar.v\#L126}{\texttt{omegabar.v:126}} --- $\bar\omega$ is invariant under digraph isomorphism.
\item[\texttt{omegabar\_sub}] (lemma) \href{https://github.com/LLM4Rocq/digraph-theory/blob/main/theories/invariants/omegabar.v\#L137}{\texttt{omegabar.v:137}} --- Sub-tournaments and vertex deletion.
\item[\texttt{omegabar\_del}] (lemma) \href{https://github.com/LLM4Rocq/digraph-theory/blob/main/theories/invariants/omegabar.v\#L143}{\texttt{omegabar.v:143}}
\item[\texttt{omegabar\_nil}] (lemma) \href{https://github.com/LLM4Rocq/digraph-theory/blob/main/theories/invariants/omegabar.v\#L148}{\texttt{omegabar.v:148}} --- $\bar\omega$ on degenerate and tiny tournaments.
\item[\texttt{omegabar\_card\_le2}] (lemma) \href{https://github.com/LLM4Rocq/digraph-theory/blob/main/theories/invariants/omegabar.v\#L156}{\texttt{omegabar.v:156}}
\item[\texttt{omegabar\_TT}] (lemma) \href{https://github.com/LLM4Rocq/digraph-theory/blob/main/theories/invariants/omegabar.v\#L164}{\texttt{omegabar.v:164}} --- ** The two reference values
\item[\texttt{omegabar\_C3}] (lemma) \href{https://github.com/LLM4Rocq/digraph-theory/blob/main/theories/invariants/omegabar.v\#L169}{\texttt{omegabar.v:169}}
\end{description}

\section*{\texttt{theories/invariants/strong.v}}
\begin{description}\setlength\itemsep{0pt}\small
\item[\texttt{strongb}] (definition) \href{https://github.com/LLM4Rocq/digraph-theory/blob/main/theories/invariants/strong.v\#L37}{\texttt{strong.v:37}}
\item[\texttt{strongP}] (lemma) \href{https://github.com/LLM4Rocq/digraph-theory/blob/main/theories/invariants/strong.v\#L39}{\texttt{strong.v:39}}
\item[\texttt{arc\_closed}] (definition) \href{https://github.com/LLM4Rocq/digraph-theory/blob/main/theories/invariants/strong.v\#L44}{\texttt{strong.v:44}} --- ** S1: forward-closed sets and the sink trick
\item[\texttt{closed\_connect}] (lemma) \href{https://github.com/LLM4Rocq/digraph-theory/blob/main/theories/invariants/strong.v\#L46}{\texttt{strong.v:46}}
\item[\texttt{rset}] (definition) \href{https://github.com/LLM4Rocq/digraph-theory/blob/main/theories/invariants/strong.v\#L54}{\texttt{strong.v:54}}
\item[\texttt{rset\_id}] (lemma) \href{https://github.com/LLM4Rocq/digraph-theory/blob/main/theories/invariants/strong.v\#L56}{\texttt{strong.v:56}}
\item[\texttt{rset\_closed}] (lemma) \href{https://github.com/LLM4Rocq/digraph-theory/blob/main/theories/invariants/strong.v\#L59}{\texttt{strong.v:59}}
\item[\texttt{rset\_trans}] (lemma) \href{https://github.com/LLM4Rocq/digraph-theory/blob/main/theories/invariants/strong.v\#L65}{\texttt{strong.v:65}}
\item[\texttt{sink\_exists}] (lemma) \href{https://github.com/LLM4Rocq/digraph-theory/blob/main/theories/invariants/strong.v\#L73}{\texttt{strong.v:73}} --- The sink trick: a reachable set of minimum size is closed and strongly connected inside.
\item[\texttt{outdeg\_closed}] (lemma) \href{https://github.com/LLM4Rocq/digraph-theory/blob/main/theories/invariants/strong.v\#L91}{\texttt{strong.v:91}} --- Closed sets keep out-degrees in the induced subgraph.
\item[\texttt{connect\_induced\_closed}] (lemma) \href{https://github.com/LLM4Rocq/digraph-theory/blob/main/theories/invariants/strong.v\#L100}{\texttt{strong.v:100}} --- Connectivity relativizes to closed sets.
\item[\texttt{connect\_cross}] (lemma) \href{https://github.com/LLM4Rocq/digraph-theory/blob/main/theories/invariants/strong.v\#L116}{\texttt{strong.v:116}} --- ** S2: cut crossing
\item[\texttt{maxpath\_no\_loopback}] (lemma) \href{https://github.com/LLM4Rocq/digraph-theory/blob/main/theories/invariants/strong.v\#L129}{\texttt{strong.v:129}} --- ** K-a1: no loopback arc at a maximum path's endpoint
\item[\texttt{disjoint\_cycles\_path}] (lemma) \href{https://github.com/LLM4Rocq/digraph-theory/blob/main/theories/invariants/strong.v\#L161}{\texttt{strong.v:161}} --- ** K-10: disjoint cycles force long paths
\item[\texttt{strongb\_induced\_closed}] (lemma) \href{https://github.com/LLM4Rocq/digraph-theory/blob/main/theories/invariants/strong.v\#L255}{\texttt{strong.v:255}} --- Strongness of the induced subgraph on a closed, internally connected set.
\item[\texttt{reduction}] (theorem) \href{https://github.com/LLM4Rocq/digraph-theory/blob/main/theories/invariants/strong.v\#L267}{\texttt{strong.v:267}} --- ** R: the composed out-regular reduction (dossier item R)
\end{description}

\section*{\texttt{theories/invariants/critical.v}}
\begin{description}\setlength\itemsep{0pt}\small
\item[\texttt{kcritical}] (definition) \href{https://github.com/LLM4Rocq/digraph-theory/blob/main/theories/invariants/critical.v\#L21}{\texttt{critical.v:21}}
\item[\texttt{kcriticalP}] (lemma) \href{https://github.com/LLM4Rocq/digraph-theory/blob/main/theories/invariants/critical.v\#L24}{\texttt{critical.v:24}}
\item[\texttt{card\_del}] (lemma) \href{https://github.com/LLM4Rocq/digraph-theory/blob/main/theories/invariants/critical.v\#L34}{\texttt{critical.v:34}} --- Deleting a vertex removes exactly one vertex.
\item[$\bigstar$ \texttt{kcritical\_proper\_sub}] (lemma) \href{https://github.com/LLM4Rocq/digraph-theory/blob/main/theories/invariants/critical.v\#L46}{\texttt{critical.v:46}} --- ** Proper subtournaments of critical tournaments In a k-$\bar\omega$-critical tournament every *proper* subtournament has $\bar\omega$ $\le$ k -- 1: it omits some vertex, hence embeds into that vertex's deletion (the Question-5.9-failure mechanism, stated once; G1 of docs/k34\_dossier.md).
\item[\texttt{C3\_kcritical2}] (lemma) \href{https://github.com/LLM4Rocq/digraph-theory/blob/main/theories/invariants/critical.v\#L64}{\texttt{critical.v:64}} --- ** C3 is 2-$\bar\omega$-critical
\item[\texttt{kcritical2\_card3}] (lemma) \href{https://github.com/LLM4Rocq/digraph-theory/blob/main/theories/invariants/critical.v\#L95}{\texttt{critical.v:95}} --- A 2-critical tournament has no 4th vertex: deleting it would leave the triangle, contradicting $\bar\omega$(T -- x) = 1.
\item[$\bigstar$ \texttt{kcritical2\_uniq}] (theorem) \href{https://github.com/LLM4Rocq/digraph-theory/blob/main/theories/invariants/critical.v\#L124}{\texttt{critical.v:124}} --- The explicit isomorphism C3 $\cong$ T sending 0,1,2 to the triangle.
\end{description}

\section*{\texttt{theories/invariants/domination.v}}
\begin{description}\setlength\itemsep{0pt}\small
\item[\texttt{dominatesb}] (definition) \href{https://github.com/LLM4Rocq/digraph-theory/blob/main/theories/invariants/domination.v\#L28}{\texttt{domination.v:28}}
\item[\texttt{dominatesbP}] (lemma) \href{https://github.com/LLM4Rocq/digraph-theory/blob/main/theories/invariants/domination.v\#L31}{\texttt{domination.v:31}}
\item[\texttt{dominatesb\_setT}] (lemma) \href{https://github.com/LLM4Rocq/digraph-theory/blob/main/theories/invariants/domination.v\#L43}{\texttt{domination.v:43}}
\item[\texttt{domnum}] (definition) \href{https://github.com/LLM4Rocq/digraph-theory/blob/main/theories/invariants/domination.v\#L46}{\texttt{domination.v:46}}
\item[\texttt{domnum\_min}] (lemma) \href{https://github.com/LLM4Rocq/digraph-theory/blob/main/theories/invariants/domination.v\#L48}{\texttt{domination.v:48}}
\item[\texttt{domnum\_witness}] (lemma) \href{https://github.com/LLM4Rocq/digraph-theory/blob/main/theories/invariants/domination.v\#L54}{\texttt{domination.v:54}}
\item[\texttt{greedy\_clique\_dom}] (lemma) \href{https://github.com/LLM4Rocq/digraph-theory/blob/main/theories/invariants/domination.v\#L63}{\texttt{domination.v:63}} --- The greedy chain: inside any \texttt{A} there is a subset that dominates \texttt{A} and is a clique of the backedge graph of \texttt{p} (the $\prec$-minimum first, every later pick beating --- and sitting $\prec$-after --- all earlier ones).
\item[$\bigstar$ \texttt{domnum\_le\_omegabar}] (theorem) \href{https://github.com/LLM4Rocq/digraph-theory/blob/main/theories/invariants/domination.v\#L114}{\texttt{domination.v:114}} --- Paper Property 3.2 --- the lower-bound engine for $\bar\omega$.
\end{description}

\section*{\texttt{theories/constructions/automorphism.v}}
\begin{description}\setlength\itemsep{0pt}\small
\item[\texttt{autb}] (definition) \href{https://github.com/LLM4Rocq/digraph-theory/blob/main/theories/constructions/automorphism.v\#L25}{\texttt{automorphism.v:25}} --- \texttt{p} is an automorphism: it preserves arcs (hence also non-arcs).
\item[\texttt{autbP}] (lemma) \href{https://github.com/LLM4Rocq/digraph-theory/blob/main/theories/constructions/automorphism.v\#L28}{\texttt{automorphism.v:28}}
\item[\texttt{dgaut}] (definition) \href{https://github.com/LLM4Rocq/digraph-theory/blob/main/theories/constructions/automorphism.v\#L35}{\texttt{automorphism.v:35}}
\item[\texttt{dgautE}] (lemma) \href{https://github.com/LLM4Rocq/digraph-theory/blob/main/theories/constructions/automorphism.v\#L37}{\texttt{automorphism.v:37}}
\item[\texttt{dgaut1}] (lemma) \href{https://github.com/LLM4Rocq/digraph-theory/blob/main/theories/constructions/automorphism.v\#L40}{\texttt{automorphism.v:40}}
\item[\texttt{dgautM}] (lemma) \href{https://github.com/LLM4Rocq/digraph-theory/blob/main/theories/constructions/automorphism.v\#L43}{\texttt{automorphism.v:43}}
\item[\texttt{dgaut\_group\_set}] (lemma) \href{https://github.com/LLM4Rocq/digraph-theory/blob/main/theories/constructions/automorphism.v\#L49}{\texttt{automorphism.v:49}}
\item[\texttt{vertex\_transitiveb}] (definition) \href{https://github.com/LLM4Rocq/digraph-theory/blob/main/theories/constructions/automorphism.v\#L55}{\texttt{automorphism.v:55}} --- Vertex-transitivity.
\item[\texttt{vertex\_transitivebP}] (lemma) \href{https://github.com/LLM4Rocq/digraph-theory/blob/main/theories/constructions/automorphism.v\#L59}{\texttt{automorphism.v:59}}
\end{description}

\section*{\texttt{theories/constructions/product.v}}
\begin{description}\setlength\itemsep{0pt}\small
\item[\texttt{lexprod}] (definition) \href{https://github.com/LLM4Rocq/digraph-theory/blob/main/theories/constructions/product.v\#L25}{\texttt{product.v:25}}
\item[\texttt{lexprod\_arcE}] (lemma) \href{https://github.com/LLM4Rocq/digraph-theory/blob/main/theories/constructions/product.v\#L31}{\texttt{product.v:31}}
\item[\texttt{card\_lexprod}] (lemma) \href{https://github.com/LLM4Rocq/digraph-theory/blob/main/theories/constructions/product.v\#L35}{\texttt{product.v:35}}
\item[\texttt{lexprod\_irrefl}] (fact) \href{https://github.com/LLM4Rocq/digraph-theory/blob/main/theories/constructions/product.v\#L45}{\texttt{product.v:45}}
\item[\texttt{lexprod\_total}] (fact) \href{https://github.com/LLM4Rocq/digraph-theory/blob/main/theories/constructions/product.v\#L50}{\texttt{product.v:50}}
\item[\texttt{lexprod\_tournament}] (definition) \href{https://github.com/LLM4Rocq/digraph-theory/blob/main/theories/constructions/product.v\#L65}{\texttt{product.v:65}} --- Object-level version for statements.
\item[\texttt{lexprod\_vertex\_transitive}] (lemma) \href{https://github.com/LLM4Rocq/digraph-theory/blob/main/theories/constructions/product.v\#L73}{\texttt{product.v:73}}
\end{description}

\section*{\texttt{theories/constructions/cayley.v}}
\begin{description}\setlength\itemsep{0pt}\small
\item[\texttt{cayley}] (definition) \href{https://github.com/LLM4Rocq/digraph-theory/blob/main/theories/constructions/cayley.v\#L30}{\texttt{cayley.v:30}} --- The carrier is the group; the (otherwise unused) connection set \texttt{A} is a real parameter of the type alias so that distinct connection sets carry distinct canonical instances.
\item[\texttt{cayley\_arcE}] (lemma) \href{https://github.com/LLM4Rocq/digraph-theory/blob/main/theories/constructions/cayley.v\#L39}{\texttt{cayley.v:39}}
\item[\texttt{cayley\_irreflP}] (lemma) \href{https://github.com/LLM4Rocq/digraph-theory/blob/main/theories/constructions/cayley.v\#L44}{\texttt{cayley.v:44}} --- ** Tournament characterization
\item[\texttt{cayley\_totalP}] (lemma) \href{https://github.com/LLM4Rocq/digraph-theory/blob/main/theories/constructions/cayley.v\#L51}{\texttt{cayley.v:51}}
\item[\texttt{translation}] (definition) \href{https://github.com/LLM4Rocq/digraph-theory/blob/main/theories/constructions/cayley.v\#L66}{\texttt{cayley.v:66}} --- ** Translations and vertex-transitivity
\item[\texttt{translationE}] (lemma) \href{https://github.com/LLM4Rocq/digraph-theory/blob/main/theories/constructions/cayley.v\#L68}{\texttt{cayley.v:68}}
\item[\texttt{translation\_aut}] (lemma) \href{https://github.com/LLM4Rocq/digraph-theory/blob/main/theories/constructions/cayley.v\#L71}{\texttt{cayley.v:71}}
\item[\texttt{cayley\_vertex\_transitive}] (lemma) \href{https://github.com/LLM4Rocq/digraph-theory/blob/main/theories/constructions/cayley.v\#L78}{\texttt{cayley.v:78}}
\end{description}

\section*{\texttt{theories/constructions/circulant.v}}
\begin{description}\setlength\itemsep{0pt}\small
\item[\texttt{Zp\_mulgE}] (lemma) \href{https://github.com/LLM4Rocq/digraph-theory/blob/main/theories/constructions/circulant.v\#L36}{\texttt{circulant.v:36}}
\item[\texttt{Zp\_invgE}] (lemma) \href{https://github.com/LLM4Rocq/digraph-theory/blob/main/theories/constructions/circulant.v\#L39}{\texttt{circulant.v:39}}
\item[\texttt{Zp\_1gE}] (lemma) \href{https://github.com/LLM4Rocq/digraph-theory/blob/main/theories/constructions/circulant.v\#L42}{\texttt{circulant.v:42}}
\item[\texttt{val\_Zp\_opp}] (lemma) \href{https://github.com/LLM4Rocq/digraph-theory/blob/main/theories/constructions/circulant.v\#L45}{\texttt{circulant.v:45}}
\item[\texttt{circulant\_arcE}] (lemma) \href{https://github.com/LLM4Rocq/digraph-theory/blob/main/theories/constructions/circulant.v\#L50}{\texttt{circulant.v:50}} --- Circulant arc characterization: \texttt{x --> y} iff the difference is a connection element.
\item[\texttt{ACset}] (definition) \href{https://github.com/LLM4Rocq/digraph-theory/blob/main/theories/constructions/circulant.v\#L63}{\texttt{circulant.v:63}}
\item[\texttt{ACset\_cond}] (lemma) \href{https://github.com/LLM4Rocq/digraph-theory/blob/main/theories/constructions/circulant.v\#L72}{\texttt{circulant.v:72}} --- The connection set hits each pair \{z, --z\} (z $\ne$ 0) exactly once --- this is the residue-arithmetic heart of "AC is a tournament".
\item[\texttt{AC\_irrefl}] (fact) \href{https://github.com/LLM4Rocq/digraph-theory/blob/main/theories/constructions/circulant.v\#L109}{\texttt{circulant.v:109}} --- AC is a tournament.
\item[\texttt{AC\_total}] (fact) \href{https://github.com/LLM4Rocq/digraph-theory/blob/main/theories/constructions/circulant.v\#L115}{\texttt{circulant.v:115}}
\item[\texttt{AC}] (definition) \href{https://github.com/LLM4Rocq/digraph-theory/blob/main/theories/constructions/circulant.v\#L124}{\texttt{circulant.v:124}}
\item[\texttt{AC\_arcE}] (lemma) \href{https://github.com/LLM4Rocq/digraph-theory/blob/main/theories/constructions/circulant.v\#L126}{\texttt{circulant.v:126}}
\item[$\bigstar$ \texttt{card\_AC}] (lemma) \href{https://github.com/LLM4Rocq/digraph-theory/blob/main/theories/constructions/circulant.v\#L129}{\texttt{circulant.v:129}}
\item[$\bigstar$ \texttt{AC\_vertex\_transitive}] (lemma) \href{https://github.com/LLM4Rocq/digraph-theory/blob/main/theories/constructions/circulant.v\#L134}{\texttt{circulant.v:134}} --- AC is vertex-transitive (as any Cayley digraph: translations act transitively) --- the M2 exit fact feeding the M3 criticality reduction.
\item[\texttt{C3\_vertex\_transitive}] (lemma) \href{https://github.com/LLM4Rocq/digraph-theory/blob/main/theories/constructions/circulant.v\#L146}{\texttt{circulant.v:146}} --- ** C is vertex-transitive \texttt{C3}'s arc relation \texttt{v == u + 1} over 'Z\_3 is translation-invariant, so the translations x  x + t are automorphisms and act transitively (the Cayley-translation argument in miniature; G3 of docs/k34\_dossier.md --- feeds the k = 4 criticality reduction).
\end{description}

\section*{\texttt{theories/invariants\_advanced/transitive.v}}
\begin{description}\setlength\itemsep{0pt}\small
\item[\texttt{del\_aut\_iso}] (lemma) \href{https://github.com/LLM4Rocq/digraph-theory/blob/main/theories/invariants_advanced/transitive.v\#L27}{\texttt{transitive.v:27}} --- An automorphism sending \texttt{u} to \texttt{g u} restricts to an isomorphism between the vertex-deleted sub-tournaments, so $\bar\omega$ agrees on them.
\item[$\bigstar$ \texttt{omegabar\_del\_vt}] (theorem) \href{https://github.com/LLM4Rocq/digraph-theory/blob/main/theories/invariants_advanced/transitive.v\#L47}{\texttt{transitive.v:47}} --- The marquee theorem: on a vertex-transitive tournament, $\bar\omega$ after deletion is the same whichever vertex is deleted.
\item[$\bigstar$ \texttt{vt\_kcritical}] (theorem) \href{https://github.com/LLM4Rocq/digraph-theory/blob/main/theories/invariants_advanced/transitive.v\#L57}{\texttt{transitive.v:57}} --- Hence k-$\bar\omega$-criticality reduces to a single deletion.
\item[\texttt{AC\_del\_uniform}] (lemma) \href{https://github.com/LLM4Rocq/digraph-theory/blob/main/theories/invariants_advanced/transitive.v\#L70}{\texttt{transitive.v:70}} --- Demo: in the circulant tournament AC all single-vertex deletions have the same $\bar\omega$ --- its criticality will need only the deletion at 0 (M4).
\end{description}

\section*{\texttt{theories/invariants\_advanced/substitution.v}}
\begin{description}\setlength\itemsep{0pt}\small
\item[\texttt{hmin}] (definition) \href{https://github.com/LLM4Rocq/digraph-theory/blob/main/theories/invariants_advanced/substitution.v\#L46}{\texttt{substitution.v:46}} --- The $\prec$-minimum of a block, and its position.
\item[\texttt{posmin}] (definition) \href{https://github.com/LLM4Rocq/digraph-theory/blob/main/theories/invariants_advanced/substitution.v\#L47}{\texttt{substitution.v:47}}
\item[\texttt{posmin\_le}] (lemma) \href{https://github.com/LLM4Rocq/digraph-theory/blob/main/theories/invariants_advanced/substitution.v\#L49}{\texttt{substitution.v:49}}
\item[\texttt{pos\_inj}] (lemma) \href{https://github.com/LLM4Rocq/digraph-theory/blob/main/theories/invariants_advanced/substitution.v\#L52}{\texttt{substitution.v:52}}
\item[\texttt{posmin\_neq}] (lemma) \href{https://github.com/LLM4Rocq/digraph-theory/blob/main/theories/invariants_advanced/substitution.v\#L55}{\texttt{substitution.v:55}}
\item[\texttt{rS\_irr}] (fact) \href{https://github.com/LLM4Rocq/digraph-theory/blob/main/theories/invariants_advanced/substitution.v\#L63}{\texttt{substitution.v:63}}
\item[\texttt{rS\_trans}] (fact) \href{https://github.com/LLM4Rocq/digraph-theory/blob/main/theories/invariants_advanced/substitution.v\#L64}{\texttt{substitution.v:64}}
\item[\texttt{rS\_total}] (fact) \href{https://github.com/LLM4Rocq/digraph-theory/blob/main/theories/invariants_advanced/substitution.v\#L65}{\texttt{substitution.v:65}}
\item[\texttt{rH\_irr}] (fact) \href{https://github.com/LLM4Rocq/digraph-theory/blob/main/theories/invariants_advanced/substitution.v\#L73}{\texttt{substitution.v:73}}
\item[\texttt{rH\_trans}] (fact) \href{https://github.com/LLM4Rocq/digraph-theory/blob/main/theories/invariants_advanced/substitution.v\#L75}{\texttt{substitution.v:75}}
\item[\texttt{rH\_total}] (fact) \href{https://github.com/LLM4Rocq/digraph-theory/blob/main/theories/invariants_advanced/substitution.v\#L77}{\texttt{substitution.v:77}}
\item[\texttt{cross\_edgeE}] (lemma) \href{https://github.com/LLM4Rocq/digraph-theory/blob/main/theories/invariants_advanced/substitution.v\#L88}{\texttt{substitution.v:88}} --- Edge correspondences between \texttt{backedge p} and the auxiliary graphs.
\item[\texttt{block\_edgeE}] (lemma) \href{https://github.com/LLM4Rocq/digraph-theory/blob/main/theories/invariants_advanced/substitution.v\#L97}{\texttt{substitution.v:97}}
\item[\texttt{omegab\_at\_lexprod\_ge}] (lemma) \href{https://github.com/LLM4Rocq/digraph-theory/blob/main/theories/invariants_advanced/substitution.v\#L103}{\texttt{substitution.v:103}} --- The bound, for this arbitrary order.
\item[$\bigstar$ \texttt{omegabar\_lexprod\_ge}] (theorem) \href{https://github.com/LLM4Rocq/digraph-theory/blob/main/theories/invariants_advanced/substitution.v\#L168}{\texttt{substitution.v:168}} --- The substitution lower bound: $\bar\omega$(S\texttt{H}) $\ge$ $\bar\omega$(S) + $\bar\omega$(H) -- 1.
\end{description}

\section*{\texttt{theories/applications/k5/acn\_arc\_facts.v}}
\begin{description}\setlength\itemsep{0pt}\small
\item[\texttt{val\_addE}] (lemma) \href{https://github.com/LLM4Rocq/digraph-theory/blob/main/theories/applications/k5/acn_arc_facts.v\#L40}{\texttt{acn\_arc\_facts.v:40}} --- ** Values of sums, opposites, differences in 'Z\_n
\item[\texttt{val\_oppE}] (lemma) \href{https://github.com/LLM4Rocq/digraph-theory/blob/main/theories/applications/k5/acn_arc_facts.v\#L43}{\texttt{acn\_arc\_facts.v:43}}
\item[\texttt{val\_sub\_le}] (lemma) \href{https://github.com/LLM4Rocq/digraph-theory/blob/main/theories/applications/k5/acn_arc_facts.v\#L46}{\texttt{acn\_arc\_facts.v:46}}
\item[\texttt{val\_sub\_gt}] (lemma) \href{https://github.com/LLM4Rocq/digraph-theory/blob/main/theories/applications/k5/acn_arc_facts.v\#L60}{\texttt{acn\_arc\_facts.v:60}}
\item[\texttt{AC\_mem\_val}] (lemma) \href{https://github.com/LLM4Rocq/digraph-theory/blob/main/theories/applications/k5/acn_arc_facts.v\#L71}{\texttt{acn\_arc\_facts.v:71}} --- ** Connection-set membership by value
\item[\texttt{AC\_arc\_lt}] (lemma) \href{https://github.com/LLM4Rocq/digraph-theory/blob/main/theories/applications/k5/acn_arc_facts.v\#L77}{\texttt{acn\_arc\_facts.v:77}} --- ** Gap characterizations under the value order
\item[\texttt{AC\_arc\_gt}] (lemma) \href{https://github.com/LLM4Rocq/digraph-theory/blob/main/theories/applications/k5/acn_arc_facts.v\#L84}{\texttt{acn\_arc\_facts.v:84}}
\item[\texttt{AC\_band\_arc}] (lemma) \href{https://github.com/LLM4Rocq/digraph-theory/blob/main/theories/applications/k5/acn_arc_facts.v\#L104}{\texttt{acn\_arc\_facts.v:104}} --- Within a band (gap < m), arcs follow the value order --- the working form of paper (iii).
\item[\texttt{AC\_mem\_Hi}] (lemma) \href{https://github.com/LLM4Rocq/digraph-theory/blob/main/theories/applications/k5/acn_arc_facts.v\#L110}{\texttt{acn\_arc\_facts.v:110}} --- ** The band facts --- paper (i) and (ii)
\item[\texttt{AC\_mem\_Hi\_opp}] (lemma) \href{https://github.com/LLM4Rocq/digraph-theory/blob/main/theories/applications/k5/acn_arc_facts.v\#L117}{\texttt{acn\_arc\_facts.v:117}}
\item[\texttt{AC\_mem\_Lo}] (lemma) \href{https://github.com/LLM4Rocq/digraph-theory/blob/main/theories/applications/k5/acn_arc_facts.v\#L138}{\texttt{acn\_arc\_facts.v:138}}
\item[\texttt{AC\_mem\_Lo\_opp}] (lemma) \href{https://github.com/LLM4Rocq/digraph-theory/blob/main/theories/applications/k5/acn_arc_facts.v\#L147}{\texttt{acn\_arc\_facts.v:147}}
\end{description}

\section*{\texttt{theories/applications/k5/acn\_base.v}}
\begin{description}\setlength\itemsep{0pt}\small
\item[\texttt{bucket\_bound}] (lemma) \href{https://github.com/LLM4Rocq/digraph-theory/blob/main/theories/applications/k5/acn_base.v\#L36}{\texttt{acn\_base.v:36}} --- ** A pigeonhole principle: sparse values in a bounded range
\item[\texttt{rid\_irr}] (fact) \href{https://github.com/LLM4Rocq/digraph-theory/blob/main/theories/applications/k5/acn_base.v\#L70}{\texttt{acn\_base.v:70}}
\item[\texttt{rid\_trans}] (fact) \href{https://github.com/LLM4Rocq/digraph-theory/blob/main/theories/applications/k5/acn_base.v\#L71}{\texttt{acn\_base.v:71}}
\item[\texttt{rid\_total}] (fact) \href{https://github.com/LLM4Rocq/digraph-theory/blob/main/theories/applications/k5/acn_base.v\#L72}{\texttt{acn\_base.v:72}}
\item[\texttt{backedge\_qid\_dist}] (lemma) \href{https://github.com/LLM4Rocq/digraph-theory/blob/main/theories/applications/k5/acn_base.v\#L79}{\texttt{acn\_base.v:79}} --- Every backedge of the value order spans a gap of at least m.
\item[\texttt{omegab\_at\_qid\_le3}] (lemma) \href{https://github.com/LLM4Rocq/digraph-theory/blob/main/theories/applications/k5/acn_base.v\#L94}{\texttt{acn\_base.v:94}} --- ** Upper bound: $\bar\omega$(AC) $\le$ 3
\item[\texttt{omegabar\_AC\_le3}] (lemma) \href{https://github.com/LLM4Rocq/digraph-theory/blob/main/theories/applications/k5/acn_base.v\#L106}{\texttt{acn\_base.v:106}}
\item[\texttt{r0\_irr}] (fact) \href{https://github.com/LLM4Rocq/digraph-theory/blob/main/theories/applications/k5/acn_base.v\#L115}{\texttt{acn\_base.v:115}}
\item[\texttt{r0\_trans}] (fact) \href{https://github.com/LLM4Rocq/digraph-theory/blob/main/theories/applications/k5/acn_base.v\#L116}{\texttt{acn\_base.v:116}}
\item[\texttt{r0\_total}] (fact) \href{https://github.com/LLM4Rocq/digraph-theory/blob/main/theories/applications/k5/acn_base.v\#L117}{\texttt{acn\_base.v:117}}
\item[\texttt{omegab\_at\_q0\_le2}] (lemma) \href{https://github.com/LLM4Rocq/digraph-theory/blob/main/theories/applications/k5/acn_base.v\#L123}{\texttt{acn\_base.v:123}}
\item[\texttt{omegabar\_ACdel\_le2}] (lemma) \href{https://github.com/LLM4Rocq/digraph-theory/blob/main/theories/applications/k5/acn_base.v\#L147}{\texttt{acn\_base.v:147}}
\item[\texttt{mem0A}] (fact) \href{https://github.com/LLM4Rocq/digraph-theory/blob/main/theories/applications/k5/acn_base.v\#L155}{\texttt{acn\_base.v:155}}
\item[\texttt{oppA\_disjoint}] (fact) \href{https://github.com/LLM4Rocq/digraph-theory/blob/main/theories/applications/k5/acn_base.v\#L158}{\texttt{acn\_base.v:158}}
\item[\texttt{cardA}] (fact) \href{https://github.com/LLM4Rocq/digraph-theory/blob/main/theories/applications/k5/acn_base.v\#L166}{\texttt{acn\_base.v:166}}
\item[\texttt{cardN0}] (fact) \href{https://github.com/LLM4Rocq/digraph-theory/blob/main/theories/applications/k5/acn_base.v\#L192}{\texttt{acn\_base.v:192}}
\item[\texttt{mem\_shift}] (fact) \href{https://github.com/LLM4Rocq/digraph-theory/blob/main/theories/applications/k5/acn_base.v\#L198}{\texttt{acn\_base.v:198}}
\item[\texttt{shift\_comp}] (fact) \href{https://github.com/LLM4Rocq/digraph-theory/blob/main/theories/applications/k5/acn_base.v\#L204}{\texttt{acn\_base.v:204}}
\item[\texttt{card\_shift}] (fact) \href{https://github.com/LLM4Rocq/digraph-theory/blob/main/theories/applications/k5/acn_base.v\#L207}{\texttt{acn\_base.v:207}}
\item[\texttt{shiftI}] (fact) \href{https://github.com/LLM4Rocq/digraph-theory/blob/main/theories/applications/k5/acn_base.v\#L210}{\texttt{acn\_base.v:210}}
\item[\texttt{shift0}] (fact) \href{https://github.com/LLM4Rocq/digraph-theory/blob/main/theories/applications/k5/acn_base.v\#L214}{\texttt{acn\_base.v:214}}
\item[\texttt{autocorr\_sym}] (fact) \href{https://github.com/LLM4Rocq/digraph-theory/blob/main/theories/applications/k5/acn_base.v\#L217}{\texttt{acn\_base.v:217}}
\item[\texttt{val1}] (fact) \href{https://github.com/LLM4Rocq/digraph-theory/blob/main/theories/applications/k5/acn_base.v\#L224}{\texttt{acn\_base.v:224}}
\item[\texttt{val\_zeta}] (fact) \href{https://github.com/LLM4Rocq/digraph-theory/blob/main/theories/applications/k5/acn_base.v\#L229}{\texttt{acn\_base.v:229}}
\item[\texttt{N0\_val\_in}] (fact) \href{https://github.com/LLM4Rocq/digraph-theory/blob/main/theories/applications/k5/acn_base.v\#L235}{\texttt{acn\_base.v:235}}
\item[\texttt{N0\_0}] (fact) \href{https://github.com/LLM4Rocq/digraph-theory/blob/main/theories/applications/k5/acn_base.v\#L244}{\texttt{acn\_base.v:244}}
\item[\texttt{autocorr\_lo}] (lemma) \href{https://github.com/LLM4Rocq/digraph-theory/blob/main/theories/applications/k5/acn_base.v\#L249}{\texttt{acn\_base.v:249}} --- The autocorrelation lemma, low half: for 0 < val t $\le$ m, the sets N and N + t share at least two elements (m $\ge$ 3).
\item[\texttt{autocorr2}] (lemma) \href{https://github.com/LLM4Rocq/digraph-theory/blob/main/theories/applications/k5/acn_base.v\#L310}{\texttt{acn\_base.v:310}} --- The autocorrelation lemma: every nonzero shift overlaps N in $\ge$ 2 points.
\item[\texttt{domnum\_AC\_ge3}] (lemma) \href{https://github.com/LLM4Rocq/digraph-theory/blob/main/theories/applications/k5/acn_base.v\#L326}{\texttt{acn\_base.v:326}} --- dom(AC) $\ge$ 3: one or two translates of N cannot cover 'Z\_n.
\item[$\bigstar$ \texttt{omegabar\_AC}] (theorem) \href{https://github.com/LLM4Rocq/digraph-theory/blob/main/theories/applications/k5/acn_base.v\#L374}{\texttt{acn\_base.v:374}} --- ** $\bar\omega$(AC) = 3
\item[\texttt{ACdel\_Ntransb}] (lemma) \href{https://github.com/LLM4Rocq/digraph-theory/blob/main/theories/applications/k5/acn_base.v\#L386}{\texttt{acn\_base.v:386}}
\item[\texttt{omegabar\_ACdel0}] (lemma) \href{https://github.com/LLM4Rocq/digraph-theory/blob/main/theories/applications/k5/acn_base.v\#L412}{\texttt{acn\_base.v:412}} --- ** $\bar\omega$(AC -- v) = 2 and 3-criticality
\item[$\bigstar$ \texttt{omegabar\_AC\_del}] (theorem) \href{https://github.com/LLM4Rocq/digraph-theory/blob/main/theories/applications/k5/acn_base.v\#L423}{\texttt{acn\_base.v:423}}
\item[$\bigstar$ \texttt{AC\_kcritical3}] (theorem) \href{https://github.com/LLM4Rocq/digraph-theory/blob/main/theories/applications/k5/acn_base.v\#L431}{\texttt{acn\_base.v:431}}
\end{description}

\section*{\texttt{theories/applications/acn\_bands.v}}
\begin{description}\setlength\itemsep{0pt}\small
\item[\texttt{radix\_ltA}] (lemma) \href{https://github.com/LLM4Rocq/digraph-theory/blob/main/theories/applications/acn_bands.v\#L34}{\texttt{acn\_bands.v:34}} --- ** Radix lemmas for lexicographic keys
\item[\texttt{radix\_lt\_inv}] (lemma) \href{https://github.com/LLM4Rocq/digraph-theory/blob/main/theories/applications/acn_bands.v\#L43}{\texttt{acn\_bands.v:43}}
\item[\texttt{radix\_eq\_inv}] (lemma) \href{https://github.com/LLM4Rocq/digraph-theory/blob/main/theories/applications/acn_bands.v\#L53}{\texttt{acn\_bands.v:53}}
\item[\texttt{band}] (definition) \href{https://github.com/LLM4Rocq/digraph-theory/blob/main/theories/applications/acn_bands.v\#L72}{\texttt{acn\_bands.v:72}} --- ** Bands
\item[\texttt{band1P}] (lemma) \href{https://github.com/LLM4Rocq/digraph-theory/blob/main/theories/applications/acn_bands.v\#L75}{\texttt{acn\_bands.v:75}}
\item[\texttt{band2P}] (lemma) \href{https://github.com/LLM4Rocq/digraph-theory/blob/main/theories/applications/acn_bands.v\#L81}{\texttt{acn\_bands.v:81}}
\item[\texttt{band3P}] (lemma) \href{https://github.com/LLM4Rocq/digraph-theory/blob/main/theories/applications/acn_bands.v\#L87}{\texttt{acn\_bands.v:87}}
\item[\texttt{band\_ge1}] (lemma) \href{https://github.com/LLM4Rocq/digraph-theory/blob/main/theories/applications/acn_bands.v\#L92}{\texttt{acn\_bands.v:92}}
\item[\texttt{band\_le3}] (lemma) \href{https://github.com/LLM4Rocq/digraph-theory/blob/main/theories/applications/acn_bands.v\#L95}{\texttt{acn\_bands.v:95}}
\item[\texttt{band12P}] (lemma) \href{https://github.com/LLM4Rocq/digraph-theory/blob/main/theories/applications/acn_bands.v\#L99}{\texttt{acn\_bands.v:99}} --- Two band dichotomies.
\item[\texttt{AC\_wrapF}] (lemma) \href{https://github.com/LLM4Rocq/digraph-theory/blob/main/theories/applications/acn_bands.v\#L114}{\texttt{acn\_bands.v:114}} --- A backward residue over a gap below m is never a connection element.
\item[\texttt{band\_gap}] (lemma) \href{https://github.com/LLM4Rocq/digraph-theory/blob/main/theories/applications/acn_bands.v\#L128}{\texttt{acn\_bands.v:128}} --- Within a (nonzero) band, value gaps stay below m.
\end{description}

\section*{\texttt{theories/applications/k5/k5\_lower.v}}
\begin{description}\setlength\itemsep{0pt}\small
\item[\texttt{ACm\_pos}] (fact) \href{https://github.com/LLM4Rocq/digraph-theory/blob/main/theories/applications/k5/k5_lower.v\#L40}{\texttt{k5\_lower.v:40}}
\item[\texttt{ACdel\_pos}] (fact) \href{https://github.com/LLM4Rocq/digraph-theory/blob/main/theories/applications/k5/k5_lower.v\#L43}{\texttt{k5\_lower.v:43}}
\item[\texttt{omegabar\_T\_ge5}] (theorem) \href{https://github.com/LLM4Rocq/digraph-theory/blob/main/theories/applications/k5/k5_lower.v\#L48}{\texttt{k5\_lower.v:48}} --- ** $\bar\omega$(T) $\ge$ 5
\item[\texttt{emb\_mem}] (fact) \href{https://github.com/LLM4Rocq/digraph-theory/blob/main/theories/applications/k5/k5_lower.v\#L62}{\texttt{k5\_lower.v:62}}
\item[\texttt{f\_inj}] (fact) \href{https://github.com/LLM4Rocq/digraph-theory/blob/main/theories/applications/k5/k5_lower.v\#L71}{\texttt{k5\_lower.v:71}}
\item[\texttt{f\_arc}] (fact) \href{https://github.com/LLM4Rocq/digraph-theory/blob/main/theories/applications/k5/k5_lower.v\#L77}{\texttt{k5\_lower.v:77}}
\item[\texttt{omegabar\_Tdel\_ge4}] (theorem) \href{https://github.com/LLM4Rocq/digraph-theory/blob/main/theories/applications/k5/k5_lower.v\#L83}{\texttt{k5\_lower.v:83}}
\end{description}

\section*{\texttt{theories/applications/k5/in\_neighbourhood.v}}
\begin{description}\setlength\itemsep{0pt}\small
\item[\texttt{gval\_bounds}] (lemma) \href{https://github.com/LLM4Rocq/digraph-theory/blob/main/theories/applications/k5/in_neighbourhood.v\#L32}{\texttt{in\_neighbourhood.v:32}} --- Connection-set elements have value in \texttt{1, m+1} $\setminus$ \{m\}.
\item[\texttt{H17\_xrange}] (lemma) \href{https://github.com/LLM4Rocq/digraph-theory/blob/main/theories/applications/k5/in_neighbourhood.v\#L42}{\texttt{in\_neighbourhood.v:42}} --- An in-neighbour of an Hi vertex sits in the window \texttt{vx--m--1, vx--1}.
\item[\texttt{H17\_dichot}] (lemma) \href{https://github.com/LLM4Rocq/digraph-theory/blob/main/theories/applications/k5/in_neighbourhood.v\#L66}{\texttt{in\_neighbourhood.v:66}} --- An in-neighbour of a Lo vertex is either below it or $\ge$ m positions up.
\item[\texttt{H17\_side}] (lemma) \href{https://github.com/LLM4Rocq/digraph-theory/blob/main/theories/applications/k5/in_neighbourhood.v\#L83}{\texttt{in\_neighbourhood.v:83}} --- H17: the common in-neighbourhood lies in the band determined by x -- y.
\item[\texttt{H17\_no\_mixed}] (lemma) \href{https://github.com/LLM4Rocq/digraph-theory/blob/main/theories/applications/k5/in_neighbourhood.v\#L107}{\texttt{in\_neighbourhood.v:107}} --- The working corollary: an Hi vertex and a Lo vertex can never both beat both an Hi vertex x and a Lo vertex y.
\end{description}

\section*{\texttt{theories/applications/k5/cells.v}}
\begin{description}\setlength\itemsep{0pt}\small
\item[\texttt{cidx}] (definition) \href{https://github.com/LLM4Rocq/digraph-theory/blob/main/theories/applications/k5/cells.v\#L45}{\texttt{cells.v:45}} --- ** Cells and the key
\item[\texttt{cidx\_decode}] (lemma) \href{https://github.com/LLM4Rocq/digraph-theory/blob/main/theories/applications/k5/cells.v\#L47}{\texttt{cells.v:47}}
\item[\texttt{cidx\_le8}] (lemma) \href{https://github.com/LLM4Rocq/digraph-theory/blob/main/theories/applications/k5/cells.v\#L56}{\texttt{cells.v:56}}
\item[\texttt{key}] (definition) \href{https://github.com/LLM4Rocq/digraph-theory/blob/main/theories/applications/k5/cells.v\#L65}{\texttt{cells.v:65}}
\item[\texttt{key\_inj}] (lemma) \href{https://github.com/LLM4Rocq/digraph-theory/blob/main/theories/applications/k5/cells.v\#L67}{\texttt{cells.v:67}}
\item[\texttt{cidx\_mono}] (lemma) \href{https://github.com/LLM4Rocq/digraph-theory/blob/main/theories/applications/k5/cells.v\#L75}{\texttt{cells.v:75}}
\item[\texttt{key\_samecell}] (lemma) \href{https://github.com/LLM4Rocq/digraph-theory/blob/main/theories/applications/k5/cells.v\#L81}{\texttt{cells.v:81}}
\item[\texttt{rk}] (definition) \href{https://github.com/LLM4Rocq/digraph-theory/blob/main/theories/applications/k5/cells.v\#L94}{\texttt{cells.v:94}} --- ** The realized key order
\item[\texttt{rk\_irr}] (fact) \href{https://github.com/LLM4Rocq/digraph-theory/blob/main/theories/applications/k5/cells.v\#L95}{\texttt{cells.v:95}}
\item[\texttt{rk\_trans}] (fact) \href{https://github.com/LLM4Rocq/digraph-theory/blob/main/theories/applications/k5/cells.v\#L96}{\texttt{cells.v:96}}
\item[\texttt{rk\_total}] (fact) \href{https://github.com/LLM4Rocq/digraph-theory/blob/main/theories/applications/k5/cells.v\#L97}{\texttt{cells.v:97}}
\item[\texttt{qk}] (definition) \href{https://github.com/LLM4Rocq/digraph-theory/blob/main/theories/applications/k5/cells.v\#L103}{\texttt{cells.v:103}}
\item[\texttt{qkE}] (lemma) \href{https://github.com/LLM4Rocq/digraph-theory/blob/main/theories/applications/k5/cells.v\#L105}{\texttt{cells.v:105}}
\item[\texttt{beat\_key}] (lemma) \href{https://github.com/LLM4Rocq/digraph-theory/blob/main/theories/applications/k5/cells.v\#L115}{\texttt{cells.v:115}}
\item[\texttt{beat\_blocks}] (lemma) \href{https://github.com/LLM4Rocq/digraph-theory/blob/main/theories/applications/k5/cells.v\#L126}{\texttt{cells.v:126}} --- Between different blocks, the beaten block-difference is a connection element; inside a block the inner difference is.
\item[\texttt{beat\_inner}] (lemma) \href{https://github.com/LLM4Rocq/digraph-theory/blob/main/theories/applications/k5/cells.v\#L136}{\texttt{cells.v:136}}
\item[\texttt{cidx\_inj\_clique}] (lemma) \href{https://github.com/LLM4Rocq/digraph-theory/blob/main/theories/applications/k5/cells.v\#L148}{\texttt{cells.v:148}} --- ** The cell Lemma: at most one clique vertex per cell
\item[\texttt{occ}] (definition) \href{https://github.com/LLM4Rocq/digraph-theory/blob/main/theories/applications/k5/cells.v\#L190}{\texttt{cells.v:190}} --- ** Occupancy and the cardinality bridge
\item[\texttt{occP}] (lemma) \href{https://github.com/LLM4Rocq/digraph-theory/blob/main/theories/applications/k5/cells.v\#L192}{\texttt{cells.v:192}}
\item[\texttt{card\_clique\_cidx}] (lemma) \href{https://github.com/LLM4Rocq/digraph-theory/blob/main/theories/applications/k5/cells.v\#L199}{\texttt{cells.v:199}}
\end{description}

\section*{\texttt{theories/applications/k5/obstructions.v}}
\begin{description}\setlength\itemsep{0pt}\small
\item[\texttt{obstrb}] (definition) \href{https://github.com/LLM4Rocq/digraph-theory/blob/main/theories/applications/k5/obstructions.v\#L30}{\texttt{obstructions.v:30}} --- The 20-disjunct obstruction pattern over the 8 cell-occupancy bits.
\item[\texttt{nat\_n\_split}] (lemma) \href{https://github.com/LLM4Rocq/digraph-theory/blob/main/theories/applications/k5/obstructions.v\#L46}{\texttt{obstructions.v:46}}
\item[\texttt{mem\_sub\_Lo\_Hi1}] (lemma) \href{https://github.com/LLM4Rocq/digraph-theory/blob/main/theories/applications/k5/obstructions.v\#L55}{\texttt{obstructions.v:55}} --- ** Membership computations for cross-band differences
\item[\texttt{mem\_sub\_Hi1\_Lo}] (lemma) \href{https://github.com/LLM4Rocq/digraph-theory/blob/main/theories/applications/k5/obstructions.v\#L72}{\texttt{obstructions.v:72}}
\item[\texttt{mem\_sub\_Hi\_m}] (lemma) \href{https://github.com/LLM4Rocq/digraph-theory/blob/main/theories/applications/k5/obstructions.v\#L90}{\texttt{obstructions.v:90}}
\item[\texttt{mem\_sub\_m\_Hi}] (lemma) \href{https://github.com/LLM4Rocq/digraph-theory/blob/main/theories/applications/k5/obstructions.v\#L107}{\texttt{obstructions.v:107}}
\item[\texttt{opp\_arcs\_contra}] (lemma) \href{https://github.com/LLM4Rocq/digraph-theory/blob/main/theories/applications/k5/obstructions.v\#L133}{\texttt{obstructions.v:133}} --- Opposite residues cannot both be connection elements.
\item[\texttt{repA\_band}] (lemma) \href{https://github.com/LLM4Rocq/digraph-theory/blob/main/theories/applications/k5/obstructions.v\#L149}{\texttt{obstructions.v:149}}
\item[\texttt{rep\_Hi}] (lemma) \href{https://github.com/LLM4Rocq/digraph-theory/blob/main/theories/applications/k5/obstructions.v\#L155}{\texttt{obstructions.v:155}}
\item[\texttt{rep\_Lo}] (lemma) \href{https://github.com/LLM4Rocq/digraph-theory/blob/main/theories/applications/k5/obstructions.v\#L162}{\texttt{obstructions.v:162}}
\item[\texttt{rep\_Z}] (lemma) \href{https://github.com/LLM4Rocq/digraph-theory/blob/main/theories/applications/k5/obstructions.v\#L169}{\texttt{obstructions.v:169}}
\item[\texttt{beatc}] (lemma) \href{https://github.com/LLM4Rocq/digraph-theory/blob/main/theories/applications/k5/obstructions.v\#L177}{\texttt{obstructions.v:177}}
\item[\texttt{bands\_neq}] (lemma) \href{https://github.com/LLM4Rocq/digraph-theory/blob/main/theories/applications/k5/obstructions.v\#L187}{\texttt{obstructions.v:187}}
\item[\texttt{triple\_A}] (lemma) \href{https://github.com/LLM4Rocq/digraph-theory/blob/main/theories/applications/k5/obstructions.v\#L200}{\texttt{obstructions.v:200}} --- ** The seven core obstruction shapes *) (* (X, zero, X) triples --- paper sets 1--8
\item[\texttt{triple\_B}] (lemma) \href{https://github.com/LLM4Rocq/digraph-theory/blob/main/theories/applications/k5/obstructions.v\#L261}{\texttt{obstructions.v:261}} --- ** The seven core obstruction shapes *) (* (X, zero, X) triples --- paper sets 1--8 *) Lemma triple\_A (u v w : backedge (qk m')) (i j k : nat) : u \textbackslash{}in K -> v \textbackslash{}in K -> w \textbackslash{}in K -> cidx (u : T5) = i -> cidx (v : T5) = j -> cidx (w : T5) = k -> (i < j)\%N -> (j < k)\%N -> (j \%\% 3)\%N = 2\%N -> ((i \%\% 3 == 0) \&\& (k \%\% 3 == 0))\%N || ((i \%\% 3 == 1) \&\& (k \%\% 3 == 1))\%N -> False.
\item[\texttt{quad\_A}] (lemma) \href{https://github.com/LLM4Rocq/digraph-theory/blob/main/theories/applications/k5/obstructions.v\#L286}{\texttt{obstructions.v:286}} --- ** The seven core obstruction shapes *) (* (X, zero, X) triples --- paper sets 1--8 *) Lemma triple\_A (u v w : backedge (qk m')) (i j k : nat) : u \textbackslash{}in K -> v \textbackslash{}in K -> w \textbackslash{}in K -> cidx (u : T5) = i -> cidx (v : T5) = j -> cidx (w : T5) = k -> (i < j)\%N -> (j < k)\%N -> (j \%\% 3)\%N = 2\%N -> ((i \%\% 3 == 0) \&\& (k \%\% 3 == 0))\%N || ((i \%\% 3 == 1) \&\& (k \%\% 3 == 1))\%N -> False.
\item[\texttt{quad\_B}] (lemma) \href{https://github.com/LLM4Rocq/digraph-theory/blob/main/theories/applications/k5/obstructions.v\#L320}{\texttt{obstructions.v:320}} --- ** The seven core obstruction shapes *) (* (X, zero, X) triples --- paper sets 1--8 *) Lemma triple\_A (u v w : backedge (qk m')) (i j k : nat) : u \textbackslash{}in K -> v \textbackslash{}in K -> w \textbackslash{}in K -> cidx (u : T5) = i -> cidx (v : T5) = j -> cidx (w : T5) = k -> (i < j)\%N -> (j < k)\%N -> (j \%\% 3)\%N = 2\%N -> ((i \%\% 3 == 0) \&\& (k \%\% 3 == 0))\%N || ((i \%\% 3 == 1) \&\& (k \%\% 3 == 1))\%N -> False.
\item[\texttt{quad\_C}] (lemma) \href{https://github.com/LLM4Rocq/digraph-theory/blob/main/theories/applications/k5/obstructions.v\#L361}{\texttt{obstructions.v:361}} --- ** The seven core obstruction shapes *) (* (X, zero, X) triples --- paper sets 1--8 *) Lemma triple\_A (u v w : backedge (qk m')) (i j k : nat) : u \textbackslash{}in K -> v \textbackslash{}in K -> w \textbackslash{}in K -> cidx (u : T5) = i -> cidx (v : T5) = j -> cidx (w : T5) = k -> (i < j)\%N -> (j < k)\%N -> (j \%\% 3)\%N = 2\%N -> ((i \%\% 3 == 0) \&\& (k \%\% 3 == 0))\%N || ((i \%\% 3 == 1) \&\& (k \%\% 3 == 1))\%N -> False.
\item[\texttt{quad\_D}] (lemma) \href{https://github.com/LLM4Rocq/digraph-theory/blob/main/theories/applications/k5/obstructions.v\#L412}{\texttt{obstructions.v:412}} --- ** The seven core obstruction shapes *) (* (X, zero, X) triples --- paper sets 1--8 *) Lemma triple\_A (u v w : backedge (qk m')) (i j k : nat) : u \textbackslash{}in K -> v \textbackslash{}in K -> w \textbackslash{}in K -> cidx (u : T5) = i -> cidx (v : T5) = j -> cidx (w : T5) = k -> (i < j)\%N -> (j < k)\%N -> (j \%\% 3)\%N = 2\%N -> ((i \%\% 3 == 0) \&\& (k \%\% 3 == 0))\%N || ((i \%\% 3 == 1) \&\& (k \%\% 3 == 1))\%N -> False.
\item[\texttt{square}] (lemma) \href{https://github.com/LLM4Rocq/digraph-theory/blob/main/theories/applications/k5/obstructions.v\#L450}{\texttt{obstructions.v:450}} --- ** The seven core obstruction shapes *) (* (X, zero, X) triples --- paper sets 1--8 *) Lemma triple\_A (u v w : backedge (qk m')) (i j k : nat) : u \textbackslash{}in K -> v \textbackslash{}in K -> w \textbackslash{}in K -> cidx (u : T5) = i -> cidx (v : T5) = j -> cidx (w : T5) = k -> (i < j)\%N -> (j < k)\%N -> (j \%\% 3)\%N = 2\%N -> ((i \%\% 3 == 0) \&\& (k \%\% 3 == 0))\%N || ((i \%\% 3 == 1) \&\& (k \%\% 3 == 1))\%N -> False.
\item[\texttt{occ\_and3F}] (lemma) \href{https://github.com/LLM4Rocq/digraph-theory/blob/main/theories/applications/k5/obstructions.v\#L495}{\texttt{obstructions.v:495}} --- ** Packaging: the occupancy bits never match an obstruction
\item[\texttt{occ\_and4F}] (lemma) \href{https://github.com/LLM4Rocq/digraph-theory/blob/main/theories/applications/k5/obstructions.v\#L505}{\texttt{obstructions.v:505}}
\item[\texttt{no\_obstruction}] (lemma) \href{https://github.com/LLM4Rocq/digraph-theory/blob/main/theories/applications/k5/obstructions.v\#L517}{\texttt{obstructions.v:517}}
\end{description}

\section*{\texttt{theories/applications/k5/coverage.v}}
\begin{description}\setlength\itemsep{0pt}\small
\item[\texttt{coverage5}] (lemma) \href{https://github.com/LLM4Rocq/digraph-theory/blob/main/theories/applications/k5/coverage.v\#L19}{\texttt{coverage.v:19}}
\end{description}

\section*{\texttt{theories/applications/k5/k5\_upper.v}}
\begin{description}\setlength\itemsep{0pt}\small
\item[\texttt{omegab\_at\_qk\_le5}] (lemma) \href{https://github.com/LLM4Rocq/digraph-theory/blob/main/theories/applications/k5/k5_upper.v\#L45}{\texttt{k5\_upper.v:45}} --- ** $\bar\omega$(T) $\le$ 5
\item[\texttt{omegabar\_T\_le5}] (theorem) \href{https://github.com/LLM4Rocq/digraph-theory/blob/main/theories/applications/k5/k5_upper.v\#L56}{\texttt{k5\_upper.v:56}}
\item[\texttt{rd\_irr}] (fact) \href{https://github.com/LLM4Rocq/digraph-theory/blob/main/theories/applications/k5/k5_upper.v\#L68}{\texttt{k5\_upper.v:68}}
\item[\texttt{rd\_trans}] (fact) \href{https://github.com/LLM4Rocq/digraph-theory/blob/main/theories/applications/k5/k5_upper.v\#L71}{\texttt{k5\_upper.v:71}}
\item[\texttt{rd\_total}] (fact) \href{https://github.com/LLM4Rocq/digraph-theory/blob/main/theories/applications/k5/k5_upper.v\#L74}{\texttt{k5\_upper.v:74}}
\item[\texttt{omegab\_at\_qd\_le4}] (lemma) \href{https://github.com/LLM4Rocq/digraph-theory/blob/main/theories/applications/k5/k5_upper.v\#L83}{\texttt{k5\_upper.v:83}}
\item[\texttt{omegabar\_Tdel\_le4}] (theorem) \href{https://github.com/LLM4Rocq/digraph-theory/blob/main/theories/applications/k5/k5_upper.v\#L109}{\texttt{k5\_upper.v:109}}
\end{description}

\section*{\texttt{theories/applications/k5/main.v}}
\begin{description}\setlength\itemsep{0pt}\small
\item[\texttt{T5\_vertex\_transitive}] (lemma) \href{https://github.com/LLM4Rocq/digraph-theory/blob/main/theories/applications/k5/main.v\#L46}{\texttt{main.v:46}} --- The product of the vertex-transitive AC with itself is vertex-transitive.
\item[$\bigstar$ \texttt{omegabar\_T5}] (theorem) \href{https://github.com/LLM4Rocq/digraph-theory/blob/main/theories/applications/k5/main.v\#L54}{\texttt{main.v:54}} --- ** $\bar\omega$(T) = 5
\item[$\bigstar$ \texttt{omegabar\_T5\_del}] (theorem) \href{https://github.com/LLM4Rocq/digraph-theory/blob/main/theories/applications/k5/main.v\#L62}{\texttt{main.v:62}} --- ** $\bar\omega$(T -- v) = 4 for every vertex v
\item[$\bigstar$ \texttt{T5\_kcritical5}] (theorem) \href{https://github.com/LLM4Rocq/digraph-theory/blob/main/theories/applications/k5/main.v\#L71}{\texttt{main.v:71}} --- ** T is 5-$\bar\omega$-critical
\item[$\bigstar$ \texttt{card\_T5}] (lemma) \href{https://github.com/LLM4Rocq/digraph-theory/blob/main/theories/applications/k5/main.v\#L79}{\texttt{main.v:79}} --- ** T has order n
\item[\texttt{proper\_sub\_omegabar\_le4}] (theorem) \href{https://github.com/LLM4Rocq/digraph-theory/blob/main/theories/applications/k5/main.v\#L84}{\texttt{main.v:84}} --- ** Every proper subtournament has $\bar\omega$ $\le$ 4
\item[$\bigstar$ \texttt{conjecture\_5\_10\_at\_k5}] (theorem) \href{https://github.com/LLM4Rocq/digraph-theory/blob/main/theories/applications/k5/main.v\#L108}{\texttt{main.v:108}} --- ** Conjecture 5.10 at k = 5: an infinite family For every N there is a 5-$\bar\omega$-critical tournament on more than N vertices.
\item[$\bigstar$ \texttt{question\_5\_9\_fails\_at\_k5}] (theorem) \href{https://github.com/LLM4Rocq/digraph-theory/blob/main/theories/applications/k5/main.v\#L127}{\texttt{main.v:127}} --- ** Question 5.9 fails at k = 5 No bound $\ell$(5) exists: for every L there is a 5-$\bar\omega$-critical tournament on more than L vertices whose ONLY subtournament with $\bar\omega$ $\ge$ 5 is the whole tournament.
\end{description}

\section*{\texttt{theories/applications/ck3/lemma7.v}}
\begin{description}\setlength\itemsep{0pt}\small
\item[\texttt{card\_ord\_ltn}] (lemma) \href{https://github.com/LLM4Rocq/digraph-theory/blob/main/theories/applications/ck3/lemma7.v\#L26}{\texttt{lemma7.v:26}} --- ** Ordinal segment cardinality (counting helper)
\item[\texttt{card\_ord\_seg}] (lemma) \href{https://github.com/LLM4Rocq/digraph-theory/blob/main/theories/applications/ck3/lemma7.v\#L37}{\texttt{lemma7.v:37}}
\item[\texttt{ckC}] (definition) \href{https://github.com/LLM4Rocq/digraph-theory/blob/main/theories/applications/ck3/lemma7.v\#L58}{\texttt{lemma7.v:58}} --- The endpoint cycle as a seq: the suffix of the path from index a.
\item[\texttt{ckCset}] (definition) \href{https://github.com/LLM4Rocq/digraph-theory/blob/main/theories/applications/ck3/lemma7.v\#L61}{\texttt{lemma7.v:61}} --- Its vertex set.
\item[\texttt{ckB}] (definition) \href{https://github.com/LLM4Rocq/digraph-theory/blob/main/theories/applications/ck3/lemma7.v\#L64}{\texttt{lemma7.v:64}} --- B = the out-neighbours of v\_\{a-1\} on the cycle.
\item[\texttt{ckS}] (definition) \href{https://github.com/LLM4Rocq/digraph-theory/blob/main/theories/applications/ck3/lemma7.v\#L68}{\texttt{lemma7.v:68}} --- S = B, the C-predecessors of B.
\item[\texttt{kernel\_full}] (theorem) \href{https://github.com/LLM4Rocq/digraph-theory/blob/main/theories/applications/ck3/lemma7.v\#L606}{\texttt{lemma7.v:606}} --- *** The package
\item[\texttt{kernel\_S}] (theorem) \href{https://github.com/LLM4Rocq/digraph-theory/blob/main/theories/applications/ck3/lemma7.v\#L631}{\texttt{lemma7.v:631}} --- The bare Lemma 7 witness.
\item[$\bigstar$ \texttt{lemma7}] (theorem) \href{https://github.com/LLM4Rocq/digraph-theory/blob/main/theories/applications/ck3/lemma7.v\#L646}{\texttt{lemma7.v:646}} --- ** Lemma 7 (Cheng--Keevash), in full generality
\item[$\bigstar$ \texttt{ck\_theorem4\_oriented}] (theorem) \href{https://github.com/LLM4Rocq/digraph-theory/blob/main/theories/applications/ck3/lemma7.v\#L684}{\texttt{lemma7.v:684}} --- ** Theorem 4 (oriented case) and the $\delta$ = 2 case of Conjecture 1
\item[$\bigstar$ \texttt{ck\_conj1\_delta2}] (corollary) \href{https://github.com/LLM4Rocq/digraph-theory/blob/main/theories/applications/ck3/lemma7.v\#L698}{\texttt{lemma7.v:698}}
\end{description}

\section*{\texttt{theories/applications/ck3/ck3\_main.v}}
\begin{description}\setlength\itemsep{0pt}\small
\item[$\bigstar$ \texttt{no\_short\_strong3}] (lemma) \href{https://github.com/LLM4Rocq/digraph-theory/blob/main/theories/applications/ck3/ck3_main.v\#L28}{\texttt{ck3\_main.v:28}} --- ** The endgame: no strong 3-outregular oriented digraph has $\ell$ $\le$ 5
\item[$\bigstar$ \texttt{ck\_conj1\_delta3}] (theorem) \href{https://github.com/LLM4Rocq/digraph-theory/blob/main/theories/applications/ck3/ck3_main.v\#L147}{\texttt{ck3\_main.v:147}} --- ** The main theorem: Conjecture 1 at $\delta$ = 3
\item[$\bigstar$ \texttt{ck\_conj1\_delta3\_path}] (corollary) \href{https://github.com/LLM4Rocq/digraph-theory/blob/main/theories/applications/ck3/ck3_main.v\#L158}{\texttt{ck3\_main.v:158}} --- Unfolded form: an explicit directed simple path with 6 arcs.
\item[\texttt{ck\_conj1\_at\_3}] (corollary) \href{https://github.com/LLM4Rocq/digraph-theory/blob/main/theories/applications/ck3/ck3_main.v\#L170}{\texttt{ck3\_main.v:170}} --- Conjecture-1-shaped alias.
\end{description}

\section*{\texttt{theories/applications/k4/k4\_lower.v}}
\begin{description}\setlength\itemsep{0pt}\small
\item[\texttt{ACm\_pos}] (fact) \href{https://github.com/LLM4Rocq/digraph-theory/blob/main/theories/applications/k4/k4_lower.v\#L44}{\texttt{k4\_lower.v:44}}
\item[\texttt{C3\_pos}] (fact) \href{https://github.com/LLM4Rocq/digraph-theory/blob/main/theories/applications/k4/k4_lower.v\#L47}{\texttt{k4\_lower.v:47}}
\item[\texttt{omegabar\_T4\_ge4}] (theorem) \href{https://github.com/LLM4Rocq/digraph-theory/blob/main/theories/applications/k4/k4_lower.v\#L52}{\texttt{k4\_lower.v:52}} --- ** $\bar\omega$(T4) $\ge$ 4 (L1)
\item[\texttt{emb\_mem}] (fact) \href{https://github.com/LLM4Rocq/digraph-theory/blob/main/theories/applications/k4/k4_lower.v\#L64}{\texttt{k4\_lower.v:64}}
\item[\texttt{f\_inj}] (fact) \href{https://github.com/LLM4Rocq/digraph-theory/blob/main/theories/applications/k4/k4_lower.v\#L69}{\texttt{k4\_lower.v:69}}
\item[\texttt{f\_arc}] (fact) \href{https://github.com/LLM4Rocq/digraph-theory/blob/main/theories/applications/k4/k4_lower.v\#L72}{\texttt{k4\_lower.v:72}}
\item[\texttt{omegabar\_T4del\_ge3}] (theorem) \href{https://github.com/LLM4Rocq/digraph-theory/blob/main/theories/applications/k4/k4_lower.v\#L77}{\texttt{k4\_lower.v:77}}
\end{description}

\section*{\texttt{theories/applications/k4/k4\_value.v}}
\begin{description}\setlength\itemsep{0pt}\small
\item[\texttt{dband}] (definition) \href{https://github.com/LLM4Rocq/digraph-theory/blob/main/theories/applications/k4/k4_value.v\#L34}{\texttt{k4\_value.v:34}} --- ** The inner band on C (h = 0 heavy, h $\ne$ 0 light)
\item[\texttt{dband2P}] (lemma) \href{https://github.com/LLM4Rocq/digraph-theory/blob/main/theories/applications/k4/k4_value.v\#L36}{\texttt{k4\_value.v:36}}
\item[\texttt{dband1P}] (lemma) \href{https://github.com/LLM4Rocq/digraph-theory/blob/main/theories/applications/k4/k4_value.v\#L39}{\texttt{k4\_value.v:39}}
\item[\texttt{dband\_ge1}] (lemma) \href{https://github.com/LLM4Rocq/digraph-theory/blob/main/theories/applications/k4/k4_value.v\#L42}{\texttt{k4\_value.v:42}}
\item[\texttt{dband\_le2}] (lemma) \href{https://github.com/LLM4Rocq/digraph-theory/blob/main/theories/applications/k4/k4_value.v\#L45}{\texttt{k4\_value.v:45}}
\item[\texttt{sidx}] (definition) \href{https://github.com/LLM4Rocq/digraph-theory/blob/main/theories/applications/k4/k4_value.v\#L60}{\texttt{k4\_value.v:60}} --- ** Sub-bands, classes, and the merged key
\item[\texttt{sidx\_decode}] (lemma) \href{https://github.com/LLM4Rocq/digraph-theory/blob/main/theories/applications/k4/k4_value.v\#L62}{\texttt{k4\_value.v:62}}
\item[\texttt{sidx\_le5}] (lemma) \href{https://github.com/LLM4Rocq/digraph-theory/blob/main/theories/applications/k4/k4_value.v\#L71}{\texttt{k4\_value.v:71}}
\item[\texttt{kappa}] (definition) \href{https://github.com/LLM4Rocq/digraph-theory/blob/main/theories/applications/k4/k4_value.v\#L80}{\texttt{k4\_value.v:80}}
\item[\texttt{kappa\_sidx}] (lemma) \href{https://github.com/LLM4Rocq/digraph-theory/blob/main/theories/applications/k4/k4_value.v\#L82}{\texttt{k4\_value.v:82}}
\item[\texttt{kv}] (definition) \href{https://github.com/LLM4Rocq/digraph-theory/blob/main/theories/applications/k4/k4_value.v\#L88}{\texttt{k4\_value.v:88}}
\item[\texttt{kv\_inj}] (lemma) \href{https://github.com/LLM4Rocq/digraph-theory/blob/main/theories/applications/k4/k4_value.v\#L90}{\texttt{k4\_value.v:90}}
\item[\texttt{kv\_kappa\_mono}] (lemma) \href{https://github.com/LLM4Rocq/digraph-theory/blob/main/theories/applications/k4/k4_value.v\#L98}{\texttt{k4\_value.v:98}}
\item[\texttt{kv\_lt\_t}] (lemma) \href{https://github.com/LLM4Rocq/digraph-theory/blob/main/theories/applications/k4/k4_value.v\#L105}{\texttt{k4\_value.v:105}}
\item[\texttt{kv\_samekappa}] (lemma) \href{https://github.com/LLM4Rocq/digraph-theory/blob/main/theories/applications/k4/k4_value.v\#L111}{\texttt{k4\_value.v:111}}
\item[\texttt{rv}] (definition) \href{https://github.com/LLM4Rocq/digraph-theory/blob/main/theories/applications/k4/k4_value.v\#L124}{\texttt{k4\_value.v:124}} --- ** The realized merged order
\item[\texttt{rv\_irr}] (fact) \href{https://github.com/LLM4Rocq/digraph-theory/blob/main/theories/applications/k4/k4_value.v\#L125}{\texttt{k4\_value.v:125}}
\item[\texttt{rv\_trans}] (fact) \href{https://github.com/LLM4Rocq/digraph-theory/blob/main/theories/applications/k4/k4_value.v\#L126}{\texttt{k4\_value.v:126}}
\item[\texttt{rv\_total}] (fact) \href{https://github.com/LLM4Rocq/digraph-theory/blob/main/theories/applications/k4/k4_value.v\#L127}{\texttt{k4\_value.v:127}}
\item[\texttt{qv}] (definition) \href{https://github.com/LLM4Rocq/digraph-theory/blob/main/theories/applications/k4/k4_value.v\#L133}{\texttt{k4\_value.v:133}}
\item[\texttt{qvE}] (lemma) \href{https://github.com/LLM4Rocq/digraph-theory/blob/main/theories/applications/k4/k4_value.v\#L135}{\texttt{k4\_value.v:135}}
\item[\texttt{n\_sub\_ge\_Sm}] (lemma) \href{https://github.com/LLM4Rocq/digraph-theory/blob/main/theories/applications/k4/k4_value.v\#L140}{\texttt{k4\_value.v:140}} --- ** Residue arithmetic helpers (pure nat, m-relative)
\item[\texttt{n\_sub\_gt\_Sm}] (lemma) \href{https://github.com/LLM4Rocq/digraph-theory/blob/main/theories/applications/k4/k4_value.v\#L143}{\texttt{k4\_value.v:143}}
\item[\texttt{n\_sub\_eq\_Sm}] (lemma) \href{https://github.com/LLM4Rocq/digraph-theory/blob/main/theories/applications/k4/k4_value.v\#L146}{\texttt{k4\_value.v:146}}
\item[\texttt{sub\_eq\_m}] (lemma) \href{https://github.com/LLM4Rocq/digraph-theory/blob/main/theories/applications/k4/k4_value.v\#L152}{\texttt{k4\_value.v:152}}
\item[\texttt{mem\_m\_F}] (lemma) \href{https://github.com/LLM4Rocq/digraph-theory/blob/main/theories/applications/k4/k4_value.v\#L159}{\texttt{k4\_value.v:159}} --- Connection-set refutations by value.
\item[\texttt{mem\_mid\_F}] (lemma) \href{https://github.com/LLM4Rocq/digraph-theory/blob/main/theories/applications/k4/k4_value.v\#L164}{\texttt{k4\_value.v:164}}
\item[\texttt{vals12}] (lemma) \href{https://github.com/LLM4Rocq/digraph-theory/blob/main/theories/applications/k4/k4_value.v\#L169}{\texttt{k4\_value.v:169}}
\item[\texttt{beat\_kv}] (lemma) \href{https://github.com/LLM4Rocq/digraph-theory/blob/main/theories/applications/k4/k4_value.v\#L187}{\texttt{k4\_value.v:187}}
\item[\texttt{beat\_blocksV}] (lemma) \href{https://github.com/LLM4Rocq/digraph-theory/blob/main/theories/applications/k4/k4_value.v\#L196}{\texttt{k4\_value.v:196}}
\item[\texttt{beat\_innerV}] (lemma) \href{https://github.com/LLM4Rocq/digraph-theory/blob/main/theories/applications/k4/k4_value.v\#L206}{\texttt{k4\_value.v:206}}
\item[\texttt{sidx\_inj\_clique}] (lemma) \href{https://github.com/LLM4Rocq/digraph-theory/blob/main/theories/applications/k4/k4_value.v\#L217}{\texttt{k4\_value.v:217}} --- ** The sub-band Lemma: at most one clique vertex per $\sigma$ (V1--V4)
\item[\texttt{occv}] (definition) \href{https://github.com/LLM4Rocq/digraph-theory/blob/main/theories/applications/k4/k4_value.v\#L263}{\texttt{k4\_value.v:263}} --- ** Occupancy
\item[\texttt{occvP}] (lemma) \href{https://github.com/LLM4Rocq/digraph-theory/blob/main/theories/applications/k4/k4_value.v\#L265}{\texttt{k4\_value.v:265}}
\item[\texttt{card\_clique\_sidx}] (lemma) \href{https://github.com/LLM4Rocq/digraph-theory/blob/main/theories/applications/k4/k4_value.v\#L272}{\texttt{k4\_value.v:272}}
\item[\texttt{occ24\_block}] (lemma) \href{https://github.com/LLM4Rocq/digraph-theory/blob/main/theories/applications/k4/k4_value.v\#L280}{\texttt{k4\_value.v:280}} --- ** V5: both $\kappa$ = 4 sub-bands occupied forces \{(m,0),(0,i)\}
\item[\texttt{occ24\_5F}] (lemma) \href{https://github.com/LLM4Rocq/digraph-theory/blob/main/theories/applications/k4/k4_value.v\#L312}{\texttt{k4\_value.v:312}}
\item[\texttt{occ24\_1F}] (lemma) \href{https://github.com/LLM4Rocq/digraph-theory/blob/main/theories/applications/k4/k4_value.v\#L340}{\texttt{k4\_value.v:340}}
\item[\texttt{occ5310\_F}] (lemma) \href{https://github.com/LLM4Rocq/digraph-theory/blob/main/theories/applications/k4/k4_value.v\#L392}{\texttt{k4\_value.v:392}} --- ** V6: the four-class chain (0,0) / high / low / $\kappa$2 dies
\item[\texttt{clique\_card\_le4}] (lemma) \href{https://github.com/LLM4Rocq/digraph-theory/blob/main/theories/applications/k4/k4_value.v\#L485}{\texttt{k4\_value.v:485}}
\item[\texttt{omegab\_at\_qv\_le4}] (lemma) \href{https://github.com/LLM4Rocq/digraph-theory/blob/main/theories/applications/k4/k4_value.v\#L500}{\texttt{k4\_value.v:500}} --- ** The value bound
\item[\texttt{omegabar\_T4\_le4}] (theorem) \href{https://github.com/LLM4Rocq/digraph-theory/blob/main/theories/applications/k4/k4_value.v\#L508}{\texttt{k4\_value.v:508}}
\item[$\bigstar$ \texttt{omegabar\_T4}] (theorem) \href{https://github.com/LLM4Rocq/digraph-theory/blob/main/theories/applications/k4/k4_value.v\#L511}{\texttt{k4\_value.v:511}}
\end{description}

\section*{\texttt{theories/applications/k4/k4\_del.v}}
\begin{description}\setlength\itemsep{0pt}\small
\item[\texttt{kd}] (definition) \href{https://github.com/LLM4Rocq/digraph-theory/blob/main/theories/applications/k4/k4_del.v\#L55}{\texttt{k4\_del.v:55}} --- ** The d\_then\_c key
\item[\texttt{kd\_inj}] (lemma) \href{https://github.com/LLM4Rocq/digraph-theory/blob/main/theories/applications/k4/k4_del.v\#L57}{\texttt{k4\_del.v:57}}
\item[\texttt{kd\_sidx\_mono}] (lemma) \href{https://github.com/LLM4Rocq/digraph-theory/blob/main/theories/applications/k4/k4_del.v\#L65}{\texttt{k4\_del.v:65}}
\item[\texttt{kd\_samesidx}] (lemma) \href{https://github.com/LLM4Rocq/digraph-theory/blob/main/theories/applications/k4/k4_del.v\#L72}{\texttt{k4\_del.v:72}}
\item[\texttt{rd\_irr}] (fact) \href{https://github.com/LLM4Rocq/digraph-theory/blob/main/theories/applications/k4/k4_del.v\#L87}{\texttt{k4\_del.v:87}}
\item[\texttt{rd\_trans}] (fact) \href{https://github.com/LLM4Rocq/digraph-theory/blob/main/theories/applications/k4/k4_del.v\#L90}{\texttt{k4\_del.v:90}}
\item[\texttt{rd\_total}] (fact) \href{https://github.com/LLM4Rocq/digraph-theory/blob/main/theories/applications/k4/k4_del.v\#L93}{\texttt{k4\_del.v:93}}
\item[\texttt{qd}] (definition) \href{https://github.com/LLM4Rocq/digraph-theory/blob/main/theories/applications/k4/k4_del.v\#L99}{\texttt{k4\_del.v:99}}
\item[\texttt{qdE}] (lemma) \href{https://github.com/LLM4Rocq/digraph-theory/blob/main/theories/applications/k4/k4_del.v\#L101}{\texttt{k4\_del.v:101}}
\item[\texttt{sidx\_del\_le4}] (lemma) \href{https://github.com/LLM4Rocq/digraph-theory/blob/main/theories/applications/k4/k4_del.v\#L105}{\texttt{k4\_del.v:105}} --- Survivors never occupy the deleted band $\sigma$ = 5 (D0).
\item[\texttt{beat\_kd}] (lemma) \href{https://github.com/LLM4Rocq/digraph-theory/blob/main/theories/applications/k4/k4_del.v\#L147}{\texttt{k4\_del.v:147}}
\item[\texttt{beat\_blocksD}] (lemma) \href{https://github.com/LLM4Rocq/digraph-theory/blob/main/theories/applications/k4/k4_del.v\#L157}{\texttt{k4\_del.v:157}}
\item[\texttt{beat\_innerD}] (lemma) \href{https://github.com/LLM4Rocq/digraph-theory/blob/main/theories/applications/k4/k4_del.v\#L167}{\texttt{k4\_del.v:167}}
\item[\texttt{didx\_inj\_clique}] (lemma) \href{https://github.com/LLM4Rocq/digraph-theory/blob/main/theories/applications/k4/k4_del.v\#L178}{\texttt{k4\_del.v:178}} --- ** D1: at most one clique vertex per band
\item[\texttt{occd}] (definition) \href{https://github.com/LLM4Rocq/digraph-theory/blob/main/theories/applications/k4/k4_del.v\#L218}{\texttt{k4\_del.v:218}} --- ** Occupancy
\item[\texttt{occdP}] (lemma) \href{https://github.com/LLM4Rocq/digraph-theory/blob/main/theories/applications/k4/k4_del.v\#L220}{\texttt{k4\_del.v:220}}
\item[\texttt{card\_clique\_didx}] (lemma) \href{https://github.com/LLM4Rocq/digraph-theory/blob/main/theories/applications/k4/k4_del.v\#L227}{\texttt{k4\_del.v:227}}
\item[\texttt{occ023\_F}] (lemma) \href{https://github.com/LLM4Rocq/digraph-theory/blob/main/theories/applications/k4/k4_del.v\#L237}{\texttt{k4\_del.v:237}} --- B1 + B3 + B4: B3 pins the B1 block at m+1 and the B4 block at s $\ge$ m+2; then (s,0) cannot beat (m+1,$\cdot$).
\item[\texttt{occ124\_F}] (lemma) \href{https://github.com/LLM4Rocq/digraph-theory/blob/main/theories/applications/k4/k4_del.v\#L299}{\texttt{k4\_del.v:299}} --- B2 + B3 + B5: B3 forces s' = m (D3f), pinning the B2 block at m; the B3-beats-B2 backedge then needs m $\in$ g.
\item[\texttt{occ0134\_F}] (lemma) \href{https://github.com/LLM4Rocq/digraph-theory/blob/main/theories/applications/k4/k4_del.v\#L356}{\texttt{k4\_del.v:356}} --- B1 + B2 + B4 + B5: the D3e core.
\item[\texttt{clique\_card\_le3}] (lemma) \href{https://github.com/LLM4Rocq/digraph-theory/blob/main/theories/applications/k4/k4_del.v\#L484}{\texttt{k4\_del.v:484}}
\item[\texttt{omegab\_at\_qd\_le3}] (lemma) \href{https://github.com/LLM4Rocq/digraph-theory/blob/main/theories/applications/k4/k4_del.v\#L500}{\texttt{k4\_del.v:500}} --- ** The deletion bound
\item[\texttt{omegabar\_T4del\_le3}] (theorem) \href{https://github.com/LLM4Rocq/digraph-theory/blob/main/theories/applications/k4/k4_del.v\#L508}{\texttt{k4\_del.v:508}}
\item[\texttt{omegabar\_T4del0}] (theorem) \href{https://github.com/LLM4Rocq/digraph-theory/blob/main/theories/applications/k4/k4_del.v\#L511}{\texttt{k4\_del.v:511}}
\end{description}

\section*{\texttt{theories/applications/k4/k4\_main.v}}
\begin{description}\setlength\itemsep{0pt}\small
\item[\texttt{T4\_vertex\_transitive}] (lemma) \href{https://github.com/LLM4Rocq/digraph-theory/blob/main/theories/applications/k4/k4_main.v\#L43}{\texttt{k4\_main.v:43}} --- The product of the vertex-transitive AC with the vertex-transitive C is vertex-transitive.
\item[$\bigstar$ \texttt{omegabar\_T4\_del}] (theorem) \href{https://github.com/LLM4Rocq/digraph-theory/blob/main/theories/applications/k4/k4_main.v\#L51}{\texttt{k4\_main.v:51}} --- ** $\bar\omega$(T4 -- v) = 3 for every vertex v
\item[$\bigstar$ \texttt{T4\_kcritical4}] (theorem) \href{https://github.com/LLM4Rocq/digraph-theory/blob/main/theories/applications/k4/k4_main.v\#L59}{\texttt{k4\_main.v:59}} --- ** T4 is 4-$\bar\omega$-critical
\item[$\bigstar$ \texttt{card\_T4}] (lemma) \href{https://github.com/LLM4Rocq/digraph-theory/blob/main/theories/applications/k4/k4_main.v\#L67}{\texttt{k4\_main.v:67}} --- ** T4 has order 3n
\item[$\bigstar$ \texttt{conjecture\_5\_10\_at\_k4}] (theorem) \href{https://github.com/LLM4Rocq/digraph-theory/blob/main/theories/applications/k4/k4_main.v\#L74}{\texttt{k4\_main.v:74}} --- ** Conjecture 5.10 at k = 4: an infinite family
\item[$\bigstar$ \texttt{question\_5\_9\_fails\_at\_k4}] (theorem) \href{https://github.com/LLM4Rocq/digraph-theory/blob/main/theories/applications/k4/k4_main.v\#L89}{\texttt{k4\_main.v:89}} --- ** Question 5.9 fails at k = 4
\end{description}

\section*{\texttt{theories/applications/unified.v}}
\begin{description}\setlength\itemsep{0pt}\small
\item[$\bigstar$ \texttt{conjecture\_5\_10\_at\_k3}] (theorem) \href{https://github.com/LLM4Rocq/digraph-theory/blob/main/theories/applications/unified.v\#L33}{\texttt{unified.v:33}} --- ** Conjecture 5.10 at k = 3: an infinite family
\item[$\bigstar$ \texttt{question\_5\_9\_fails\_at\_k3}] (theorem) \href{https://github.com/LLM4Rocq/digraph-theory/blob/main/theories/applications/unified.v\#L47}{\texttt{unified.v:47}} --- ** Question 5.9 fails at k = 3
\item[$\bigstar$ \texttt{conjecture\_5\_10\_at\_345}] (theorem) \href{https://github.com/LLM4Rocq/digraph-theory/blob/main/theories/applications/unified.v\#L66}{\texttt{unified.v:66}} --- ** The unified headline: every k $\in$ \{3, 4, 5\}
\end{description}

